% 本文件由 LaTeX 排版 Agent 根据有效 translation.json 自动生成。
% author=Tertullian; work_id=0317; title=驳斥\普拉克塞阿斯

\chapter{驳斥\Person{普拉克塞阿斯}}

% structure: p0001 source=h1 target=omitted reason=page-title-already-rendered-by-parent

在此书中，他在所有要点上捍卫了天主圣三的道理。\par

% structure: p0003 source=h2 target=section reason=relative-heading-rank; base=h2

\section{第一章 撒但对真理的诡计。它们如何采取\Person{普拉克塞阿斯}的异端的形式。此异端发表的记述。}

魔鬼曾以各种方式与真理竞争并抵抗真理。有时他的目的是通过捍卫真理来毁灭真理。他主张只有一位主，即世界全能的创造者，以便从这个\Emph{一体的}道理中制造出异端。他说圣父亲自降入童贞女，亲自由她而生，亲自受苦，甚至亲自是耶稣基督。在这里，\Emph{古}蛇与自己矛盾，因为当他在若翰的洗礼后试探基督时，他以\Quote{天主之子；}的身份接近祂，这肯定暗示了天主有一个儿子，甚至根据圣经的见证，他当时正从那些圣经中编制他的试探：\BibleQuote{你若是天主子，就命这些石头变成饼罢。}\BibleRef{玛 4:3}又说：\BibleQuote{你若是天主子，就从这里跳下去；因为经上记着：祂必为你吩咐祂的天使}——无疑是指圣父——\BibleQuote{他们要用手托着你，免得你的脚碰在石头上。}又或许，他只是在用谎言责备福音，实际上说：\Quote{废除玛窦！废除路加！\Emph{为什么要理会他们的话}？尽管他们这样说，\Emph{我断言}我所接近的是天主自己；我面对面试探的是全能者自己；我接近祂无别的目的，就是为了试探祂。相反，如果那\Emph{仅仅是}天主之子，我大概绝不会屈尊与祂打交道。}然而，他从起初就是撒谎者，\BibleRef{若 8:44}他按自己的方式煽动任何人；例如，\Person{普拉克塞阿斯}。因为\Emph{他}是第一个从\Place{亚细亚}把这种异端的邪恶带到\Place{罗马}的人，一个在其他方面性情不安定的人，尤其是因宣信者的骄傲而膨胀，唯独因为他曾短暂忍受监狱的烦恼；在那时候，甚至\BibleQuote{即使他把自己的身体交给火焚烧，也对他毫无益处，}没有天主的爱，\BibleRef{格前 13:3}天主的恩赐正是他所抵抗和摧毁的。因为罗马主教曾承认\Person{孟他努}、\Person{普黎斯卡}和\Person{马克西米拉}的预言恩赐，并因此将他的平安赐予\Place{亚细亚}和\Place{夫黎基雅}的教会，\Emph{他}却强行对先知们和他们的教会提出虚假指控，坚持主教的前任在教座上的权威，迫使罗马主教召回他已发出的和平信函，并停止承认\Emph{上述}恩赐。这样，\Person{普拉克塞阿斯}在罗马为魔鬼做了双重服务：他驱逐了预言，引进了异端；他赶走了圣神，并钉死了圣父。此外，\Person{普拉克塞阿斯}的莠子已经撒下，并且在此地也结了果实，当时许多人在教义的单纯中沉睡；但这些莠子似乎已被拔除，被天主乐意使用的那个人发现并揭露了。确实，\Emph{\Person{普拉克塞阿斯}}曾故意恢复他旧有的（真）信仰，在放弃错误后教导它；并且有他自己的笔迹作为证据，留存在属肉性的人中间，那次交易发生在他们的会众中；此后就没有再听到他的消息。我们这边，后来因承认和持守圣神，从属肉性的人中退出了。但是\Emph{\Person{普拉克塞阿斯}的}莠子在当时到处撒下了种子，这些种子隐藏了一段时间，其活力隐藏在面具下，现在又重新爆发了生命。但是，它必再次被连根拔起，如果主愿意，即使是现在；但如果不是现在，就在所有莠子被捆成捆的日子，连同一切绊脚石一起，用不灭的火焚烧。\BibleRef{玛 13:30}\par

% structure: p0005 source=h2 target=section reason=relative-heading-rank; base=h2

\section{第二章 天主圣三与合一的公教道理，有时称为神圣经纶，或天主位格关系的安排。}

然后，随着时间推移，圣父确实降生，圣父受苦，天主自己，全能的主，他们在宣讲中宣告祂就是耶稣基督。然而我们，正如我们一直所做的（并且我们自从被护慰者更好地教导以来，祂的确引导人进入一切真理），相信只有一个天主，但按照以下安排，或希腊文称为\OriginalTerm{安排}{\GreekText{οἰκονομία}}，即这唯一的天主也有一个圣子，祂的圣言，从祂本身发出，万物藉祂而造，没有祂什么也不能造。\Emph{我们相信}祂曾被圣父派遣到童贞女那里，并由她所生——既是人又是天主，人子与天主之子，并被称为耶稣基督之名；\Emph{我们相信}祂曾受苦、死亡并被埋葬，按照圣经，并且在被圣父复活并带回天堂后，坐在圣父的右边，\Emph{并且}祂将审判活人与死人；祂也按照自己的应许，从天上由圣父派遣了圣神，护慰者，信仰（相信圣父、圣子及圣神者）的圣化者。这信仰的规范从福音之初就传给我们，甚至在那些较古老的异端者之前，更在昨日的\Emph{冒充者}\Person{普拉克塞亚斯}之前，这一点从所有异端所标榜的晚近日期，以及我们这位新奇的普拉克塞亚斯的绝对新颖性来看，将是显而易见的。在此原则中，我们此后也必须找到一条同等有力的推定来反对一切异端——即凡是最初的便是真的，而日期较晚的则是假的。但保持这一规约规则不变，仍然必须给予机会审查（异端者的陈述），以便教导和保护各类人；至少使每项对\Emph{真理}的歪曲不会显得未经审查就被定罪，仅仅被预判；尤其是针对这一异端，它自以为拥有纯正真理，认为除了说圣父、圣子及圣神是同一位格之外，别无他法相信唯一的天主。仿佛如此\Quote{一}就不是\Quote{万有}，因为万有来自\Quote{一}，藉着（即实体的）统一；而安排的奥秘仍被保守，它将一体分配为天主圣三，将三位——圣父、圣子及圣神——按次序放置：然而，三者不是在状态上，而是在等级上；不是在实体上，而是在形式上；不是在能力上，而是在方面上；然而属于同一实体、同一状态、同一能力，因为祂是唯一的天主，这些等级、形式和方面都从祂而来，并以圣父、圣子及圣神的名号计算。它们如何能有数目而不分离，本论文将在后续证明。\par

% structure: p0007 source=h2 target=section reason=relative-heading-rank; base=h2

\section{第三章 种种常见的恐惧与偏见。天主圣三在统一中的教义从这些误解中得救}

那单纯的——我不愿称他们为愚笨或无知的——人，他们总是构成信徒的大多数，因着（天主圣三的）安排而惊愕，因为他们的信德准则正是使他们远离世上众多的神明，归向唯一真实的天主；他们不明白，虽然祂是唯一的天主，但必须按照祂自己的\OriginalTerm{安排}{\GreekText{οἰκονομία}}来相信祂。他们将天主圣三的数目顺序和分布视为对唯一性的分割；然而，那从自身中衍生出天主圣三的唯一性，非但未被摧毁，反而被其所支持。他们不断反对我们，说我们宣讲两位神、三位神，而他们自己却特别以敬拜独一神自居；就好像唯一性本身不会因不合理的推论而产生异端，而合理论证的天主圣三则构成真理。他们说：我们维护\OriginalTerm{唯一统治}{\GreekText{μοναρχία}}。于是，就发音而论，连拉丁人（且是无知者）也如此发音，你会以为他们对\OriginalTerm{唯一统治}{\GreekText{μοναρχία}}的理解与他们的发音一样完整。那么，拉丁人费力地发出\OriginalTerm{唯一统治}{\GreekText{μοναρχία}}的音，而希腊人却拒绝理解\OriginalTerm{安排}{\GreekText{οἰκονομία}}。至于我自己，如果我掌握了两种语言的任何知识，我确信\OriginalTerm{唯一统治}{\GreekText{μοναρχία}}除开单一和个体的统治外别无他意；但尽管如此，这唯一统治并不因为是一人的治理，就排除那治理者拥有一个儿子，或使自己实际成为自己的儿子，或任凭他所愿的代理人管理他自己的唯一统治。不仅如此，我主张，任何主权都不像他自己的那样只属于一人，或在如此意义上独一，或在如此意义上是一统，以至于不能通过与其最密切联结、且本身已提供作为其官员的其他人物来管理。此外，如果那拥有唯一统治的人有一个儿子，且这儿子也被接纳为分享者，这唯一统治并不会立即分裂而不再是一统；就其起源而言，它同样属于那将它传给儿子的人；既然是属于他的，它仍然是一统，因为由两个如此不可分割者维系在一起。因此，既然神的唯一统治也由如此多的天使军团和天军管理，正如经上所记：\Quote{有万万服事祂，有亿亿侍立在祂面前}\BibleRef{达 7:10}；既然它并未因由如此众多的势力管理而不再是独一的统治，那么，为何在圣子和圣神——他们被赋予第二和第三位，且在本体上与圣父如此紧密相连——的身上，天主被认为遭受分裂和割裂，而在如此众多的天使身上却未遭受任何分裂和割裂？你真的认为那些自然属于圣父本体的一部分、祂爱的保证、祂大能的工具，甚至祂的权能本身和祂整个唯一统治体系的那几位，是唯一统治的倾覆和毁灭吗？你这样想是不对的。我宁愿你锻炼自己理解事物的含义，而不是纠结于字音。现在你必须明白，唯一统治的倾覆是指这样：当另一个具有自身框架和状态的统治被引入之上并成为竞争者时；例如，如同马西翁的观点中，引入另一位神反对造物主；或如同你们的瓦伦提诺和普罗迪科那样，引入许多神。那便成了唯一统治的倾覆，因为它导致造物主的毁灭。\par

% structure: p0009 source=h2 target=section reason=relative-heading-rank; base=h2

\section{第四章 天主性体的合一与至上统治和独一治理。公教教义并未损伤唯一统治。}

但至于我，我从没有其他来源，只从圣父的本体引申圣子，并（表达祂）在圣父的旨意之外不做任何事，并已从圣父领受了一切权能，我在圣子身上保存了独一统治，正如圣父所交付给祂的一样，我怎么可能正在破坏信仰中的独一统治呢？我在谈到\Emph{天主性中的第三位}时也正式做出同样说明，因为我相信圣神从没有其他来源，只从圣父通过圣子而\Emph{发出}。那么你们要当心，破坏独一统治的恐怕是\Emph{你们}，当你们推翻了它的安排和分施，而天主乐意以如此多的名号来\Emph{使用}时。然而，尽管引入了天主圣三，独一统治在其自身状态中仍然坚定稳固，以致圣子实际上必须将它完整地归还给圣父；甚至正如宗徒在其书信中关于万物的终结所说：\Quote{当祂把王国交付给天主，即圣父时；因为祂必须为王，直到祂把一切仇敌放在祂脚下；}\BibleRef{格前 15:24}当然接着圣咏的话：\Quote{坐在我的右边，直到我使你的仇敌作你的脚凳。}\Quote{然而，当一切事物都屈服于祂时——那使一切屈服于祂的除外——那时圣子自己也要屈服于那使一切屈服于祂的，好叫天主成为万物中的万有。}\BibleRef{格前 15:27}我们因此看到圣子并不妨碍独一统治，尽管它现在由圣子管理；因为与圣子同在时，它仍处于自身状态，并将由圣子以自身状态归还给圣父。因此，没有人会因为接纳圣子而损害它，因为肯定它已由圣父托付给祂，并且将来还要由祂归还给圣父。现在，我们从\Emph{受默感的}宗徒的书信这一节中，已经能够表明圣父和圣子是两位\Emph{不同的位格}，不仅因为提到祂们不同的名字：圣父和圣子，还因为那交付王国的那一位和承受王国的那一位——同样，那使一切屈服的那一位和那一切向之屈服的那一位——必然是两位不同的存在。\par

% structure: p0011 source=h2 target=section reason=relative-heading-rank; base=h2

\section{第五章 天主的圣言或圣子由圣父以神圣的发出而生成。以人的思想和意识的活动加以说明。}

但既有人主张两者实为一位，以致圣父与圣子被视为同一，则理应详察关于圣子的全部问题：祂是否存在，祂是谁，以及祂存在的方式。如此，真理本身必藉圣经及其守护性的解释获得确认。有人声称\WorkTitle{创世纪}的希伯来文本开头为：\Quote{在起初，天主为自己造了一个儿子。}既无根据，我遂转而依据天主自身的安排——祂在创造世界之前，直至圣子受生之前所存在的状态。因为在万有之先，天主是孤独的——祂在自己内且为自己就是宇宙、空间和万有。然而，祂是孤独的，因为除祂自身外，没有外在之物。但即使那时祂也不是孤独的；因为祂拥有那在祂自身内的东西，即祂自己的\OriginalTerm{理性}{Ratio}。因为天主是理性的，理性最先在祂内；所以万有皆出自祂自身。这理性就是祂自己的思想（或意识），希腊人称之为\OriginalTerm{圣言}{\GreekText{λόγος}}，我们也用这个词表示话语或言谈；因此，如今在我们中间，由于对一词的简单解释，常有人说圣言在起初就与天主同在；不过，将理性视为更古老或许更为恰当；因为天主并非起初就有圣言，而是在起初之前就已拥有理性；因为圣言本身也由理性构成，从而证明理性是先在的，作为自身的实体。但这区别并无实际意义。因为尽管天主尚未发出祂的圣言，祂仍在自己内拥有祂，与祂的理性同在并包含在其中，正如祂默默策划和安排一切日后要藉圣言发出的事物。当祂这样与自己的理性策划安排时，祂实际上使正在以话语方式处理的东西成为圣言。为使你更易理解，首先从你自身考虑，你是按\BibleQuote{天主的肖像和模样}受造的\BibleRef{创 1:26}，为何你自身也拥有理性？你是理性的受造物，不仅由理性的工匠造成，而且实际由其本体所赋予生命。请观察，当你默默与自己交谈时，这过程就在你内藉理性进行：在你思想的每一次活动、每一念头的冲动中，理性都以话语迎接你。凡你所思，便有一个话语；凡你所想，便有理性。你必须在心中说出来；当你说话时，你承认言语是与你对话者，其中便含有这理性，你藉此在思想中与你的话语交谈，同时（经由互动）藉与话语的交谈产生思想。如此，在一定意义上，话语是你内的第二个\Quote{位格}，你通过它思考、发言，也通过它（互动地）发言、产生思想。话语本身是与你不同的东西。在天主内，这一切何等更圆满地实现！——你尚且被视为祂的肖像和模样，因为祂在沉默时也在自身内拥有理性，并且理性内含着祂的话语。故我可毫不冒昧地首先确立（固定原则）：即使在宇宙受造之前，天主也不是孤独的，因为祂在自己内拥有理性，以及内含于理性的圣言——祂藉着在自己内推动圣言，使圣言成为次于祂的。\par

% structure: p0013 source=h2 target=section reason=relative-heading-rank; base=h2

\section{第六章 天主的圣言也是天主的智慧。智慧按神圣计划出来创造宇宙}

这神圣智能的能力和安排，在圣经中也以\OriginalTerm{智慧}{\GreekText{Σοφία}}之名显明；因为有什么比天主的理性或圣言更有资格称为智慧呢？请听智慧自身——以第二位格的身份说话：\BibleQuote{在起初，主造了我作为祂道路的开始，为了祂的化工；在祂造地以前，在山岳安放以前，在丘陵受孕以前，祂已生了我}\BibleRef{箴 8:22}，也就是说，祂在自己的理智中创造并生了我。然后，请注意智慧与主相伴所暗示的区别。智慧说：\BibleQuote{当祂预备诸天时，我在那里；当祂在风——即天上的云——之上设立祂的堡垒时；当祂稳固泉源和天下一切时，我就在祂旁边，与祂一同安排一切；我就在祂所喜悦的那位旁边，每日在祂面前欢乐}\BibleRef{箴 8:27}。当天主乐意将祂在自己内与智慧理性及圣言所策划安排的事物赋予各自的实体和形式时，祂首先发出圣言本身，圣言内含有祂不可分离的理性与智慧，以便万有都藉着祂而造成——万有本已在天主的理智和智慧中被策划、安排并已造成，只是尚需公开显现并永久保存于各自的形式和实体中。\par

% structure: p0015 source=h2 target=section reason=relative-heading-rank; base=h2

\section{第七章 圣子被称为圣言和智慧（按人类思想和言语的不足），易被视为纯粹属性。祂被证明是位格性的存在}

因此，那时圣言也亲自取了祂自己的形体和荣耀的衣袍，祂自己的声音和言语，当天主说：\BibleQuote{有光！}的时候。\BibleRef{创 1:3} 这就是圣言的完美诞生，当祂从天主那里发出时——首先由祂形成，以智慧的名义构思和思考\Emph{一切}万物——\BibleQuote{主创造了（或形成了）我，作为祂道路的起始；}\BibleRef{箴 8:22} 然后被\Emph{生}，以执行一切——\Quote{当祂预备诸天时，我与祂同在。} 这样，祂使祂与自己平等：因为祂从自己发出，就成为了祂的首生子，因为在万有之先受生；\BibleRef{哥 1:15} 也是祂的独生子，因为唯独从天主而生，以祂特有的方式，从祂自己内心的腹中——正如圣父亲自作证：\Quote{我的心，}祂说，\Quote{已发出了我最卓越的圣言。}圣父永远喜悦祂，祂也同样以互惠的喜乐在圣父面前喜乐：\Quote{祢是我的圣子，我今天生了祢；}甚至在晨星之前我就生了祢。圣子也同样以智慧的名号以祂自己的身份承认圣父：\Quote{主形成了我，作为祂道路的起始，为祂自己的工程；在众山之前，祂生了我。} 因为如果智慧在这段经文中似乎说，她为主所造是为了祂的工程，并成就祂的道路，但另一处圣经却证明：\BibleQuote{万物是由圣言造成的，没有祂，什么都没有造成；}\BibleRef{若 1:3} 同样，在另一处（它说），\Quote{藉祂的圣言，诸天被建立，其中的万军由祂的圣神}——也就是说，藉着在圣言内的圣神（或天主性）：\Emph{这样}显而易见的是，同一个能力，在一处被描述为智慧的名号，在另一处被描述为圣言的称谓，祂被初创是为了天主的工程\BibleRef{箴 8:22}，祂\Quote{坚固了诸天}；\BibleQuote{万物由祂而造成}\BibleRef{若 1:3}，\BibleQuote{没有祂，什么也没有造成。}\BibleRef{若 1:3} 我们无需再在此点上多费笔墨，仿佛被提及的并不是圣言本身，祂在智慧、理性以及整个神圣灵魂与圣神的名号下被谈论。祂也成为天主子，并在祂从祂发出时受生。那么（你问），你承认圣言是由圣神和智慧的交流所构成的一种实体吗？我当然承认。但你不允许祂真正是一个有实体的存在，拥有自己的实体；这样祂可以被视为一个客观事物和一位位格，因而能够（作为位于天主圣父之后的第二位）构成两位：圣父和圣子，天主和圣言。因为你会说，言语是什么？不过是口中的声音和音响，并且（如语法学家所教）当空气被撞击时，可被耳朵理解，但其余则是一种空洞、虚无、无形体之物。我相反主张，从天主那里不会发出任何空洞虚无的东西，因为祂不是从空洞虚无中发出的；也不可能从如此伟大的实体中发出的东西缺乏实体，并且产生了如此伟大的实体：因为凡藉祂而造的一切，祂亲自（位格性地）造了。怎能说祂自己是虚无，而没有祂什么也没有造呢？怎能说空\Emph{无}的能造出坚实之物，虚空的能造出充满之物，无形体的能造出有形体之物呢？因为尽管有时被造物可能与造它的不同，但虚空空洞之物不能造出任何东西。那么，被称为圣子、本身被称为天主的圣言，难道是虚空空洞之物吗？\BibleQuote{圣言与天主同在，圣言就是天主。}\BibleRef{若 1:1} 经上记着：\BibleQuote{你不可妄呼天主的名。}\BibleRef{出 20:7} 这肯定是祂，\BibleQuote{祂虽具有天主的形体，却没有以自己与天主同等为抢夺的。}\BibleRef{斐 2:6} 具有天主的什么形体？当然是指某种形体，而不是没有形体。因为谁能否认天主是一个身体，尽管\BibleQuote{天主是神}？\BibleRef{若 4:24} 因为神有它自己种类的形体实体，有它自己的形式。现在，即使无形之物，无论是什么，在天主内都有其实体和形式，因此只有天主才看得见它们，那么从祂实体中发出的东西，岂能没有实体呢？因此，无论我所称谓的圣言的实体是什么，我为其要求圣子的名号；当我承认圣子时，我断言祂与圣父的区别，是第二位。\par

% structure: p0017 source=h2 target=section reason=relative-heading-rank; base=h2

\section{第八章　虽然天主子或圣言由圣父发出，祂却不似 \Person{瓦伦提努斯}的流溢说那样与圣父分离。圣神也不与两者分离。来自自然的说明。}

如果有人因此认为我是在引入某种\OriginalTerm{生发}{\GreekText{προβολή}}——即一物由另一物发出的过程，正如\Person{瓦伦提诺}所做的那样，他一连串地设立一个又一个\OriginalTerm{移涌}{Æon}——那么我首先要回答你：真理不应因异端也使用这一术语而回避使用它及其所表达的真实与意义。事实上，异端反倒是从真理那里借用了它，为要将其塑造成自己的仿冒品。天主的圣言究竟有没有被发出？请你在此与我站定，不要退缩。若祂被发出，那么就承认真正的教义包含一种生发；不必在意异端在某些方面模仿真理。问题在于，双方是在什么意义上使用某个事物及其表达该事物的语词。瓦伦提诺将他的种种发出与其本源分割、分离，并将它们置于距本源如此遥远之处，以至于移涌不认识父：它确实渴望认识父，却不能；不仅如此，它几乎被吞没并消解于其余的质料中。但在我们这里，只有圣子认识父\BibleRef{玛 11:27}，并且祂亲自揭示了\BibleQuote{父的怀抱}\BibleRef{若 1:18}。祂也曾与父同听同见一切事；凡父所命令祂的，祂就照样讲论\BibleRef{若 8:26}。祂所成就的不是自己的旨意，而是父的旨意\BibleRef{若 6:38}，这旨意祂从起初就极亲密地知晓。\BibleQuote{除了天主内住的神灵外，谁能知道天主的事呢？}\BibleRef{格前 2:11}。然而，圣言是由圣神形成的，并且（容我这样说）圣神是圣言的躯体。因此，圣言既常在与父内，正如祂所说：\BibleQuote{我在父内}\BibleRef{若 14:11}；又常与天主同在，正如经上所载：\BibleQuote{圣言与天主同在}\BibleRef{若 1:1}；并且从未与父分离，或不同于父，因为\BibleQuote{我与父原是一体}\BibleRef{若 10:30}。这就是由真理——合一的守护者——所教导的生发，其中我们宣称圣子是从父发出的，却未与父分离。因为天主发出了圣言，正如\OriginalTerm{护慰者}{Paraclete}所宣告的，正如树根发出树木，泉源发出河流，太阳发出光芒。这些就是\OriginalTerm{发出}{\GreekText{προβολαί}}，即从它们所由出的本体而发出的\Emph{流出物}。确实，我不妨称树木为树根的儿子或后裔，称河流为泉源的后裔，称光芒为太阳的后裔；因为每一个原始源头都是父母，凡从源头涌出的都是后裔。更不用说天主的圣言了，祂实际上获得了\Emph{圣子}这一专属于祂的称号。然而，树木并未与树根分离，河流也未与泉源分离，光芒也未与太阳分离；同样，圣言也未与天主分离。因此，遵循这些类比的形态，我承认我称天主和祂的圣言——父与祂的子——为\Emph{两位}。因为树根与树木显然是两样事物，却是相互联结的；泉源与河流也是两种形态，却不可分割；同样，太阳与光芒是两种形态，却是连贯的。凡从某物发出的，必然次于其所从出之物，但并不因此分离。然而，凡有第二者，必有两样；凡有第三者，必有三样。如今，圣神确实是从天主和圣子而来的第三位；正如树木的果实是从树根而来的第三物，或从河流流出的溪涧是从泉源而来的第三物，或光芒的尖端是从太阳而来的第三物。然而，无一物与其所由获得自身属性的原始本源相异。同样，天主圣三从父通过缠绕连接的过程流下，丝毫不扰乱\OriginalTerm{独一君权}{Monarchy}，同时守护着\OriginalTerm{经世之序}{Economy}。\par

% structure: p0019 source=h2 target=section reason=relative-heading-rank; base=h2

\section{第九章 天主教信仰准则中的一些要点阐述。尤其关于至圣天主圣三各不同位格的无混淆区分。}

你当常铭记：这就是我所宣认的信仰准则；以此准则，我证明圣父、圣子、圣神彼此不可分离，如此你便知道此话是在何种意义上说的。如今，请注意，我的断言是：圣父是独一的，圣子是独一的，圣神是独一的，并且他们彼此有别。这句话被每个无学识的人以及每个心术不正的人错误地理解为似乎在主张一种多样性，其含义暗示圣父、圣子、圣神之间是分离的。此外，当他们（在颂扬\Emph{\OriginalTerm{独一主宰}{Monarchy}}的同时贬低\Emph{\OriginalTerm{经纶}{Economy}}）主张圣父、圣子、圣神是同一的时，我不得不这样说：圣子与圣父的不同并非出于多样性，而是出于分施；他的不同不是由于分割，而是由于区别；因为圣父并不等同于圣子，他们在存有的方式上彼此有别。因为圣父是整个本体，而圣子是整体的流出和一部分，正如他自己所承认的：\BibleQuote{我的父比我大。} \BibleRef{若 14:28}在圣咏中，他的卑微被描述为\BibleQuote{稍微逊于天使。}因此，圣父与圣子有别，他大于圣子，因为生者是一位，受生者是另一位；差遣者是一位，被差遣者是另一位；创造者是一位，那藉以创造者又是另一位。幸而主亲自使用了关于护慰者位格的这种说法，以表示不是分割或分离，而是一种安排（天主性内的相互关系）；因为他说：\BibleQuote{我要祈求父，他必赐给你们另一位护慰者……就是真理之神} \BibleRef{若 14:16}，从而使护慰者与自己有别，正如我们说圣子也与圣父有别一样；因此他在护慰者身上显示了第三级，而我们相信第二级在圣子身上，这是由于\Emph{经纶}中所观察到的秩序。此外，他们拥有\Emph{圣父}和\Emph{圣子}这些不同的名号这一事实，不正等于宣告他们在位格上是有别的吗？因为，当然，一切事物都将如其名号所代表的那样；它们是什么，将永远是什么，就会被那样称呼；名号所指示的区别绝不容许任何混淆，因为在它们所指的事物中并无混淆。\BibleQuote{是就是是，非就是非；此外多余的，便是出于邪恶。} \BibleRef{玛 5:37}\par

% structure: p0021 source=h2 target=section reason=relative-heading-rank; base=h2

\section{第十章 圣父与圣子的名号本身便证明两者的位格区别。它们不可能同一，亦无需同一以保存天主的独一主宰}

所以，或者是圣父，或者是圣子；白天不同于黑夜；圣父与圣子也不相同，以至于两者应是同一，而此一或彼一应是两者——这是最为自负的\Quote{\OriginalTerm{唯一神论者}{Monarchians}}所坚持的意见。他们说，祂自己使自己成为自己的儿子。然而，圣父生圣子，圣子生圣父；那些这样相互从对方而成为彼此相关者，绝不可能独自简单地如此与自己相关，使圣父能使自己成为自己的儿子，使圣子能使自己成为自己的父亲。而天主所建立的关系，祂也守护它们。父亲必须有一个儿子，才能成为父亲；同样，儿子必须有一个父亲，才能成为儿子。然而，拥有是一回事，而存在是另一回事。例如，为了成为丈夫，我必须有一个妻子；我绝不能自己成为自己的妻子。同样，为了成为父亲，我有一个儿子，因为我绝不能是儿子对我自己；而为了成为儿子，我有一个父亲，因为我不可能是我自己的父亲。正是这些关系造就了我（我所是），当我拥有它们时：我将是父亲，当我有一个儿子时；将是儿子，当我有一个父亲时。如果我要对自己是这些关系中的任何一种，我便不再有我所是的东西：既不是父亲，因为我将是我自己的父亲；也不是儿子，因为我将是我自己的儿子。此外，既然我应当拥有这些关系中的一个，以便成为另一个；所以，如果我要同时成为两者，我将因不拥有另一个而无法成为一个。因为如果我自己必须是我的儿子，而我又是一个父亲，那么我就不再有儿子，因为我是我自己的儿子。但由于没有儿子（因为我是我自己的儿子），我怎能是父亲？因为我应当有一个儿子，才能成为父亲。因此我不是儿子，因为我没有一个父亲（是儿子之所以成为儿子者）。同样，如果我自己是我的父亲，而我又是一个儿子，那么我就不再有父亲，而是我自己是我的父亲。但由于没有父亲（因为我是我自己的父亲），我怎能是儿子？因为我应当有一个父亲，才能成为儿子。因此我不能是父亲，因为我没有一个儿子（是父亲之所以成为父亲者）。这一切必定是魔鬼的诡计——这种将一方排除和分离于另一方的做法——因为假装\Emph{\OriginalTerm{唯一神权}{Monarchy}}将两者包含在一体中，他使两者都不被持有和承认，因此祂不是圣父，因为祂确实没有圣子；也不是圣子，因为同样祂没有圣父：因为当祂是圣父时，祂不会是圣子。就这样，他们坚守\Emph{\OriginalTerm{唯一神权}{Monarchy}}，却既没有圣父，也没有圣子。然而，\BibleQuote{在天主没有不可能的事。}\BibleRef{玛 19:26} \Emph{确实如此}；谁不知道呢？又有谁不知道\BibleQuote{在人不可能的，在天主却可能}？\BibleRef{路 18:27}\BibleQuote{天主也拣选了世上愚拙的，为使那有智慧的羞愧。}\BibleRef{格前 1:27} 我们都读到过。因此，他们争辩说，对天主而言，使自己同时成为圣父和圣子并不困难，违背人间事物的状况。因为不生育的妇人违反自然生子，对天主并非难事；童贞女怀孕也是如此。当然，没有什么是\BibleQuote{为主太难的事。}\BibleRef{创 18:14} 但是，如果我们选择在我们任意的想象中如此放肆而苛刻地应用这个原则，我们就可以借此得出天主做了任何我们喜欢的事，理由是对祂来说不是不可能的。然而，我们不能因为祂能做一切事，就假定祂实际做了祂没有做的事。但我们应当询问\Emph{祂是否确实如此做了}。天主如果愿意，可以给人装上翅膀飞翔，就如祂给了鸢翅膀一样。然而，我们不能就得出结论说祂这样做了，因为祂能够这样做。祂也可以立即消灭普拉克西亚斯和所有其他异端者；然而，不能仅仅因为祂有能力，就推论祂这样做了。因为必须既有鸢也有异端者；也必须使圣父被钉十字架。在某种意义上，甚至对天主而言也有难事——即祂没有做过的事——不是因为祂不能，而是因为祂不愿意做。因为对天主而言，愿意就是能够，不愿意就是不能够；然而，凡祂所愿意的，祂既能完成，也显示了祂的能力。因此，既然如果天主愿意使自己成为自己的儿子，祂有能力这样做；既然如果祂有能力，祂就实现了\Emph{祂的目的}，那么，当你向我们证明了祂确实这样做了之后，你才能证实祂的能力和祂的意愿（甚至做这事）。\par

% structure: p0023 source=h2 target=section reason=relative-heading-rank; base=h2

\section{第十一章  \Person{普拉克西亚斯}所主张的圣父与圣子的同一性充满困惑和荒谬。许多圣经经文引用证明天主圣三各个位格的区分}

然而，你有责任像我们一样，从圣经中清楚地提出你的证据，正如我们证明天主使自己成为圣言为自己的圣子。因为如果祂称祂为圣子，而圣子无非是从\Emph{圣父}本身所发出的那一位，并且圣言是从\Emph{圣父}本身所发出的，那么祂就是圣子，而不是祂所自出的那位。因为\Emph{圣父}本身并没有从自己发出。现在，你说圣父与圣子相同，实际上就是使同一位格既从自身发出（同时又从自身出来）那位是神的存在。如果祂有可能这样做，祂至少没有这样做。你必须提出我要求你的证据——像我的证据一样；也就是说，（你必须向我证明）圣经显示圣子和圣父是相同的，正如我们这边证明圣父和圣子是有区别的；我说的是\Emph{有区别}，但\Emph{不可分离}：因为正如我提出天主自己的话：\Quote{我的心涌出了我的最优越的圣言}，同样，你也应当提出一些经文来反对我，其中天主说过：\Quote{我的心涌出了我自己作为我的最优越的圣言}，这样祂既是发出者又是被发出者，既是派遣者又是被派遣者，因为祂既是圣言又是神。我也请你注意，在我这一边，我提出了圣父对圣子所说的话：\Quote{你是我的儿子，我今日生了你}。如果你要我相信祂既是圣父又是圣子，请给我指出另一段经文，其中宣告：\Quote{主对自己说：我是我自己的儿子，我今日生了我自己}；或者又说：\Quote{在早晨之前，我生了我自己}；同样，\Quote{我，主，为我的工作占有了我自己作为我道路的开始；在众山之前，我也生了我自己}；以及其他任何类似的经文。此外，万有的主天主，如果事实如此，为什么祂会犹豫不这样对自己说呢？祂是害怕如果不这样说，就不会被相信吗？祂至少害怕一件事——说谎。祂也害怕自己，害怕自己的真理。因此，既然我相信祂是真神，我确信祂所宣告的一切都是按照祂自己的安排和计划存在的，并且祂所安排的一切都是按照祂自己的宣告。而你呢，你必须把祂说成是说谎者、欺骗者和篡改祂话语的人，如果当祂自己是自己的圣子时，祂却把圣子的角色指派给另一位，而全部圣经都清楚地证明天主圣三的存在和\Emph{位格}的区别，并且确实为我们提供了我们的\OriginalTerm{信德准则}{Rule of faith}，即说话的那位，祂所说的那位，以及祂对之说话的那位，不可能似乎是同一位。如此荒谬和误导的说法不配于天主，即当祂是对自己说话时，祂却对另一位说话，而不是对祂自己。那么，请听圣父通过依撒意亚的口关于圣子的其他话语：\Quote{看，我的儿子，我所拣选的；我所爱的，我喜悦的；我要把我的神放在祂身上，祂要向外邦人施行审判。} \BibleRef{依 42:1} 也请听祂对圣子所说的话：\Quote{你被称为我的儿子，去复兴雅各伯的支派，领回以色列分散的人，这难道是小事情吗？我立你作外邦人的光明，使我的救恩直到地极。} \BibleRef{依 49:6} 现在也请听圣子关于圣父的话语：\Quote{上主的神在我身上，因为祂用油傅了我，叫向人们传报福音。} 祂也在圣咏中对圣父这样说：\Quote{求你不要抛弃我，直到我向后来的世代宣扬你臂膀的能力。} 同样，在另一篇圣咏中：\Quote{上主，迫害我的人何其众多！} 但几乎所有预言基督位格的圣咏，都\Emph{描述}基督向天主说话。也请注意圣神以第三位的身份论及圣父和圣子的话：\Quote{主对我主说：你坐在我右边，等我使你的仇敌作你的脚凳。} 同样，在依撒意亚的话中：\Quote{上主对我的受傅者主这样说。} \BibleRef{依 45:1} 同样，在同一位先知中，祂论及圣子对圣父说：\Quote{主，谁相信了我们的报道？上主的手臂向谁显示了呢？我们关于他所传的，好像他是一个幼童，好像干地中的根苗，没有形态，也没有英华。} \BibleRef{依 53:1} 这些是从许多见证中选取的少数几个；因为我们并不打算提出全部圣经段落，因为我们的主题在各条目中已经积累了相当多的经文，正如我们在各章中将其作为具有完全尊严和权威的见证者那样。但在这少数引文中，天主圣三中\Emph{位格}的区别清楚地展现出来。因为有圣神自己说话，有圣父对之说话，有圣子被论及。同样，其他段落也确立了每一位格的独特特征——有些是对圣父或对圣子论及圣子，有些是对圣子或对圣父论及圣父，还有一些是对圣神说的。\par

% structure: p0025 source=h2 target=section reason=relative-heading-rank; base=h2

\section{第十二章 其他引自圣经的引文，证明天主内位格的多元}

如果天主圣三的数目也冒犯你，好像它与单纯的一体不相关联，我问你，一个仅仅是并且绝对是单一和独一的存有，怎么可能用复数说法说话，说：\BibleQuote{让我们照我们的肖像，按我们的样式造人；}\BibleRef{创 1:26} 而祂本该说：\Quote{让我照我的肖像，按我的样式造人，} 作为一个独特且单一的存有？然而，在接下来的经文中：\BibleQuote{看，人已变得与我们相似，}\BibleRef{创 3:22} 如果祂是独一且单一的，祂要么是在欺骗我们，要么是在戏弄我们，用复数说话。或许祂是对天使说话，正如犹太人解释这段经文那样，因为他们也不认识圣子？或者，是因为祂同时是圣父、圣子和圣神，所以用复数称呼对自己说话，因此而使自己成为复数？不，是因为祂已将祂的圣子紧贴身边，作为第二位格、祂自己的圣言，并且还有第三位格——圣言中的圣神，所以祂故意采用复数短语：\Quote{让\Emph{我们}造；} 以及\Quote{照\Emph{我们}的肖像；} 以及\Quote{变得像我们\Emph{中的一位}。} 因为祂与谁造了人？又使人与谁相似？（答案是）一方面圣子，祂将来要穿上人性；另一方面圣神，祂要圣化人。祂正是与这些，在\OriginalTerm{天主圣三}{Trinity}的一体中，如同与祂的臣仆和见证者那样说话。在接下来的经文中，祂也区分了位格：\BibleQuote{于是天主照自己的肖像造了人；造了祂，照天主的肖像。}\BibleRef{创 1:27} 为什么要说\BibleQuote{天主的肖像}？为什么不说仅仅是\BibleQuote{祂自己的肖像}，如果创造者只有一位，并且没有另一位按其肖像祂造了人？但有一位，天主按其肖像正在造人，也就是说，基督的肖像，祂将来要成为人（更确实、更真实地），已经使那人被称为祂的肖像，那人当时正要用泥土塑造——那位真实而完全的人的肖像和模样。然而，关于世界之前的创造，圣经说了什么？它的第一个陈述确实作出，当时圣子尚未显现：\BibleQuote{天主说：有光，就有了光。}\BibleRef{创 1:3} 立刻显现出圣言，\BibleQuote{那真光，照亮进入世界的每个人，}\BibleRef{若 1:9} 借着祂，光也来到了世界。从那一刻，天主愿意创造在圣言中实现，基督临在并服事祂：于是天主创造了。天主说：\BibleQuote{要有穹苍，……天主就造了穹苍；}\BibleRef{创 1:6} 天主又说：\Quote{要有光体（在穹苍中）；\Emph{于是}天主造了大光和小光。} 但其余一切受造之物，祂也照样造了，那造了前者的——我是指天主的圣言，\BibleQuote{万物是借着祂造的，凡受造的，没有一样不是借着祂造的。}\BibleRef{若 1:3} 现在，如果祂也是天主，按照若望的话（他说）\BibleQuote{圣言是天主，}\BibleRef{若 1:1} 那么你就有了两个存有——一个命令事情被造，另一个\Emph{执行命令并}创造。然而，你应当理解祂是另一位，在何种意义上，我已经解释过了，基于位格，而非实体——以区分的方式，而非分裂的方式。但我必须到处持守三个连贯不可分的（位格）中的唯一实体，然而我不得不承认，鉴于情况的必要性，发命令的那一位与执行的那一位是不同的。因为，如果祂一直在自己做工，同时命令第二位去做，祂就不会发命令了。但祂确实发了命令，尽管如果祂只是独一的一位，祂就不会打算命令自己；否则，祂一定会在没有命令的情况下工作，因为祂不会等着命令自己。\par

% structure: p0027 source=h2 target=section reason=relative-heading-rank; base=h2

\section{第十三章 各类圣经段落在位格复数性与实体一体性方面的力量示例。此处无多神论，因为一体性被强调为对抗多神论的补救。}

那么，你回答说，如果说话的是天主，创造的是天主，那么，一位天主说话，另一位创造；这样，就说了两位天主。如果你如此鲁莽和苛刻，请思考片刻；为使你更妥善、更审慎地思考，请听那篇描述两位为天主的圣咏：\Quote{天主，你的宝座永远长存；你王国的权杖是正义的权杖。你喜爱正义，憎恨邪恶：因此天主，你的天主，用喜油膏了你，或者使你成为他的基督。} 既然他在这里对天主说话，并肯定天主被天主所膏，那么，他必然因权杖的王权而肯定两位是天主。同样，依撒意亚也对基督说：\Quote{舍巴人，身材高大的人，要归顺你；他们要被锁链捆绑跟随你；他们要崇拜你，因为天主在你内：因为你是我们的天主，我们却不认识你；你是以色列的天主。} 因为在这里，通过说\Quote{天主在你内}和\Quote{你是天主}，他指出了两位天主：（在前一句\Quote{在你内}中，他指在基督内；在另一句中，他指圣神）。还有更伟大的陈述，你会在福音中明确找到：\Quote{在起初已有圣言，圣言与天主同在，圣言就是天主。} \BibleRef{若 1:1} 有一位\Quote{存在着}，另一位\Quote{与他同在}。但我在圣经中也看到\Quote{主}这个名称适用于两者：\Quote{上主对我主说：你坐在我右边。} 依撒意亚也说：\Quote{主，谁相信了我们的报道？上主的臂膀向谁显示了呢？} \BibleRef{依 53:1} 他肯定会说\Quote{你的臂膀}，如果他不希望我们理解父是主，子也是主。在创世记中，我们还有更古老的见证：\Quote{那时，上主从天上从主那里降下硫磺与火，焚烧了索多玛和哈摩辣。} \BibleRef{创 19:24} 现在，要么否认这是圣经；要么（我问你）你是怎样的人，竟然不认为词语应按其写下的意义来接受和理解，尤其是当它们不是用比喻和寓言表达，而是用确定和简单的陈述？如果你追随那些当时不能忍受主显示自己是天主之子的人，因为他们不相信他是主，那么（我问你）请和他们一起记起那经文：\Quote{我曾说过：你们是神，你们都是至高者的儿子；} 还有，\Quote{天主站在诸神集会中；} 这样，如果圣经不害怕称那些因信成为天主之子的人为神，你就可以确信，同一部圣经更恰当地将主的名赐给了天主真实而独一的儿子。很好！你说，我将从今天起（并且凭这些圣经的权威）按照你的观点宣讲两位天主和两位主。千万不要（这是我的回答）。因为我们——藉天主的恩宠拥有对圣经的时间和场合的洞察力，尤其是我们这些追随护慰者而非人间师傅的人——确实明确地宣告：\Emph{两位}存在是天主，父和子，再加上圣神，甚至\Emph{三位}，按照\OriginalTerm{安排（经济）}{Economy}的原则，它引入了\Emph{数目}，以免父如你荒谬推论的那样被认为自己诞生和受苦，这是不允许相信的，因为并未如此传授。然而，说两位天主或两位主，这话从不从我们口中说出：并非因为父是天主、子是天主、圣神是天主、每一位都是天主的说法不真实，而是因为早期确实有两位被称为天主，两位被称为主，以便当基督来时，他可以被承认为天主并被称为主，他是那既是天主又是主者的儿子。如果圣经中只找到一位天主和主的人格，那么基督就完全不能被接纳为天主和主的称号：因为（在圣经中）除了独一的天主和独一的主，没有别位，那么父就必须亲自降临（人间），因为圣经中只读到一位天主和一位主，而他的整个\OriginalTerm{安排（经济）}{Economy}就会晦暗不明，这安排（经济）在他眷顾的安排中带着如此清晰的预见被筹划和布置，作为我们信德的对象。然而，当基督来了，被我们认出他就是那从一开始就引起（神圣安排中的）多样性者，是来自父的\OriginalTerm{第二位}{second}，与圣神一起是\OriginalTerm{第三位}{third}，并且他更充分地向我们启示和显示了父（比以往任何时候），那位既是天主又是主的名号立即被恢复至（天主性体的）\OriginalTerm{合一}{Unity}，因为外邦人必须从他们的众多偶像转向唯一的天主，以便明确区分敬拜独一天主的人和信奉多神的人。因为基督徒应当作为\Quote{光明之子}照耀世界，敬拜并呼求那独一天主和主为\Quote{世界之光}。此外，如果我们从那种完美认识——即肯定天主和主的称号适用于父、子和圣神——去呼求\OriginalTerm{多位}{a plurality of}神和主，我们会熄灭自己的火炬，会变得不那么勇敢地忍受殉道的苦难，而处处都有容易逃脱的途径，只要我们像某些坚持多于一位神的异端那样，指着\OriginalTerm{多位}{a plurality of}神和主发誓。因此，我根本不会谈神，也不会谈主，而会跟随宗徒；这样，如果父和子同样受呼求，我就称父为\Quote{\Emph{天主}}，并呼求耶稣基督为\Quote{\Emph{主}}。\BibleRef{罗 1:7} 但如果只提基督，我便能称他为\Quote{\Emph{天主}}，如同同一宗徒所说：\Quote{基督是在万有之上，永远受赞美的天主。} \BibleRef{罗 9:5} 因为我也会将\Quote{\Emph{太阳}}的名称赋予一道阳光本身；但如果我提到发出光线的太阳，我当然会立刻从单纯的光线中撤回太阳的名称。虽然我不造出两个太阳，但我仍把太阳和它的光线视为同一不可分实体的两个事物和两种形式，正如天主和他的圣言，父和子。\par

% structure: p0029 source=h2 target=section reason=relative-heading-rank; base=h2

\section{第十四章。圣父的本性不可见性，以及圣子在旧约许多段落中被见证的可见性。由此提供的他们区别的论据。}

此外，当我们坚持圣父和圣子是\Emph{二位}时，那规定天主为不可见的治理原则，便来援助我们。当梅瑟在埃及渴望看见主的面容，说：\BibleQuote{因此，如果我在祢眼中得蒙恩宠，求祢将自己显示给我，使我能看见祢并认识祢}，\BibleRef{出 33:13} 天主说：\BibleQuote{你不能看见我的面容；因为没有人看见我还能存活。}换句话说，看见我的必要死。如今我们发现，天主曾被许多人看见，然而看见祂的人没有一个因此而死。\Emph{事实是}，他们是按照人的能力看见天主，而非按照天主性的圆满荣耀。因为圣祖们（如亚巴郎和雅各伯）以及先知们（如依撒意亚和厄则克耳）都被记载看见过天主，但他们并没有死。那么，他们本应死去，既然他们看见了祂——因为（经文说）\Quote{没有人能看见天主而存活}；否则，如果他们看见了天主却没有死，那么圣经就是假的，因为天主说\Quote{若有人看见我的面容，他必不能存活}。无论如何，圣经误导了我们，它既说天主不可见，又显示祂让我们看见。那么，被看见的祂必定是另一位，因为对于被看见的那位，不能说祂是不可见的。因此，结论是：以不可见者，我们当理解为在祂尊威圆满中的圣父；而我们承认圣子因祂受生之存在的安排，是可见的。就如我们不得在太阳的全部实体——即在天上的太阳——中注视它，而只能以眼睛承受一束光线，因为从它投射到地上的这束光线是经过调节的。此处，对方可能有人会争辩说：圣子作为圣言和圣神也是不可见的；并且，因主张圣父和圣子同一性，而断言圣父与圣子实为同一\Emph{位格}。但圣经，如我们所说，通过可见与不可见的区别，维护了它们的差异。他们于是进一步论证说：如果当时是圣子向梅瑟说话，祂必定是论及自己说祂的面容无人可见，因为祂自己正是以圣子之名显现的不可见的圣父。通过这些，他们想要证明可见与不可见是同一的，正如圣父与圣子是同一的；（他们如此主张）因为在前面的段落中，在祂拒绝向梅瑟显示祂的面容之前，圣经告诉我们\BibleQuote{上主与梅瑟面对面说话，如同人与朋友交谈}；\BibleRef{出 33:11} 正如雅各伯也说：\BibleQuote{我面对面见了天主。}\BibleRef{创 32:30} 因此，可见与不可见是同一的；两者如此相同，那么祂作为圣父是不可见的，作为圣子则是可见的。仿佛按照我们的解释，圣经不适用于圣子，因为圣父在祂自己的不可见性中被搁置一边。我们却宣告：圣子本身（作为圣子）也是不可见的，因为祂是天主，是天主圣言和圣神；但在祂肉身的\Emph{日子}之前，祂是可见的，祂这样对亚郎和米黎盎说：\BibleQuote{如果你们中有先知，我将在异象中向他显示自己，在梦中与他说话；但不像梅瑟，我要与他面对面说话，\Emph{明明地}，即真实地，而非\Emph{隐晦地}，}即借着形象；\BibleRef{户 12:6} 正如宗徒也表达：\BibleQuote{现在我们透过镜子模糊地（或隐晦地）看见，但那时将面对面。}\BibleRef{格前 13:12} 既然祂将祂亲自显现并与梅瑟面对面说话保留到将来——这应许后来在（显圣容的）山上实现，正如我们在福音中读到\BibleQuote{梅瑟出现，与耶稣谈话}，——那么很明显，在早期，天主，即天主子，总是借着（如）镜子和谜语、异象和梦，向先知和圣祖，也向梅瑟自己显现。而且，即使上主可能曾与梅瑟面对面说话，他也不能作为人看见祂的面容，除非是在（如）镜子和谜语中。此外，如果上主如此与梅瑟说话，以致梅瑟实际看见了祂的面容，眼对眼，那么为什么紧接着在同一场合，他又渴望看见祂的面容——这他不该渴望，因为他已经看见了？同样，上主也说祂的面容不能被看见，因为祂已经显示了它（如果祂确实显示了，如我们的对手所认为的）。或者，那被拒绝看见的天主面容是什么，如果有一个面容是人可见的？\Quote{我见了天主，}雅各伯说，\Quote{面对面，我的生命得保存。}\BibleRef{创 22:30} 应该另有某个面容，若被看见就会杀死人。那么，圣子可见吗？（当然不，）尽管祂是天主的面容，唯独在异象和梦中，在镜子和谜语中，因为（天主的）圣言和圣神不能在想象的形式之外被看见。但是，（他们说）祂称不可见的圣父为祂的面容。圣父是谁？祂不正是圣子的面容吗，因祂作为圣父所生的那一位所获得的权威？因为说某位更大的人物：\Quote{那人是我的面容；他给我他的面容}，不是有一种自然的恰当性吗？\Quote{我的父，}基督说，\Quote{比我大。}\BibleRef{若 14:28} 因此圣父必是圣子的面容。因为圣经说什么？\BibleQuote{祂位格的圣神是主基督。}因此，正如基督是圣父位格的圣神，有充分的理由，由于一体性，那祂属于其位格的那一位——即圣父——的圣神，称祂为自己的\Quote{面容}。这当然是令人惊讶的事，即圣父可以被视为圣子的面容，而祂是圣子的头；因为\BibleQuote{基督的头是天主。}\BibleRef{格前 11:3}\par

% structure: p0031 source=h2 target=section reason=relative-heading-rank; base=h2

\section{第十五章。引用新约经文。它们证实同一真理：圣子的可见性与圣父的不可见性相对照。}

倘若我未能从旧约中可能引起争议的段落来解决（我们信仰的）这一条，我将从新约中拿出支持我们观点的证据，免得你立即把一切（关系和状况）都归于圣父，而我则将它们归于圣子。看，我发现在福音书和宗徒（著作）中，揭示了一位可见的和不可见的天主，二者在状况上有明确而位格的区别。若望有一句强调的话：\BibleQuote{从来没有人看见过天主} \BibleRef{若 1:18}；当然是指以前的任何时间。但他通过说天主从未被看见，确实消除了所有时间的问题。宗徒证实了这句话；因为论到天主，他说：\BibleQuote{没有人看见过，也不能看见} \BibleRef{弟前 6:16}；因为看见祂的人必定会死。但正是这同一些宗徒证明他们曾看见并\Quote{摸过}基督 \BibleRef{若一 1:1}。那么，如果基督本身就是圣父和圣子，祂怎么可能是可见又不可见的呢？然而，为了调和可见与不可见之间的这种差异，反对者中会不会有人辩解说，这两句话完全正确：祂在肉身中确实是可见的，但在肉身显现之前是不可见的；因此，那在肉身之前作为圣父不可见的，与那在肉身中作为圣子可见的是同一位？但如果祂就是降生成人前不可见的那一位，那么祂在古时（在肉身到来之前）怎么实际上被人看见呢？按同样的推理，如果祂就是（在肉身到来之后）可见的那一位，那么宗徒们现在怎么宣告祂是不可见的呢？\Emph{这一切怎能成立}，除非那在古时只以预象和谜语的方式可见、并通过降生成人而更清晰地显现的是一位，即那成了肉身的圣言；而那从来没有人看见过的则是另一位，即圣父，就是圣言所属的那一位？总之，让我们查考宗徒们看见的是谁。若望说：\BibleQuote{那从起初就有的，我们亲眼看见过，并亲手摸过的生命之言} \BibleRef{若一 1:1}。生命之言成了肉身，被人听见、看见并摸过，因为祂是肉身，而在（祂来到）肉身之前，祂是\BibleQuote{在起初与天主同在的圣言}，与圣父同在 \BibleRef{若 1:1}，而不是圣父与圣言同在。因为圣言固然是天主，但祂与天主同在，因为祂是天主的天主；祂与圣父结合，便与圣父同在。\BibleQuote{我们见过祂的光荣，正如父独生者的光荣} \BibleRef{若 1:14}；当然是指圣子的光荣，就是那位可见的、并被不可见的圣父所光荣的。因此，既然他说过天主的圣言是天主，为了不给敌人留下（以为）他看见了圣父本身的妄断，并为了区分不可见的圣父和可见的圣子，他好像\OriginalTerm{额外丰富地}{ex abundanti}加了一句：\BibleQuote{从来没有人看见过天主} \BibleRef{若一 4:12}。他指的是哪位天主？圣言吗？但他已经说过：\BibleQuote{\Emph{祂}我们见过、听过，并亲手摸过生命之言}。那么，（我再问）他指的是哪位天主？当然是圣父，圣言与祂同在，即那独生子，祂在父怀里，并亲自向我们启示了祂 \BibleRef{若 1:18}。祂被听见、看见，并且为了不致被认为是幻影，还被摸过。保禄也见过祂；但他并没有看见圣父。他说：\BibleQuote{我不是见过耶稣\Emph{我们的主基督}吗？} \BibleRef{格前 9:1}。此外，他明确地称基督为天主，说：\BibleQuote{按血统说，基督也是从他们来的，祂是在万有之上，世世代代应受赞美的天主} \BibleRef{罗 9:5}。他也向我们指明，天主子，即天主的圣言，是可见的，因为那成了肉身的被称为基督。但论到圣父，他对弟茂德说：\BibleQuote{没有人见过，也不能看见}；他用更丰富的描述累积说：\BibleQuote{只有祂是不朽的，住在不可接近的光中} \BibleRef{弟前 6:16}。也是在前面的一节经文中，他曾论到祂说：\BibleQuote{愿尊荣归于万世的君王，那不朽坏、不可见、唯一的天主} \BibleRef{弟前 1:17}；这样我们甚至可以把相反的特质归于圣子本身——可朽坏、可接近——宗徒证明祂\BibleQuote{照经上所记载的死了} \BibleRef{格前 15:3}，并且\BibleQuote{最后也显现给我}——当然是通过可接近的光，尽管他经历那光并非没有危及视力 \BibleRef{宗 22:11}。\Emph{类似的危险也曾降临到}伯多禄、若望和雅各伯身上（他们面对同样的光）并未丧失理智；如果他们连圣子的荣耀都承受不了，而看见了圣父，他们当场就会死：\BibleQuote{因为没有人看见了我还能活着} \BibleRef{出 33:20}。既然如此，显然那位在终结时变得可见的，从起初就一直被看见；而那位从起初就从未可见的，在终结时也不被看见；因此有两位——可见的和不可见的。所以，被看见的始终是圣子，与世人交谈的始终是圣子，凭圣父的权威和旨意行事的一直是圣子；因为\BibleQuote{子不能凭自己做什么，祂看见父做什么，才能做什么} \BibleRef{若 5:19} ——\Quote{做什么}，即是在祂的意念和思想中。因为父凭意念和思想行事；而圣子，祂在父的意念和思想中，将祂所看见的予以实现和成形。因此万物是藉着圣子造的，没有祂，什么也没有造 \BibleRef{若 1:3}。\par

% structure: p0033 source=h2 target=section reason=relative-heading-rank; base=h2

\section{第十六章 天主圣子在旧约中的早期显现，以及其后降生成人的预演}

但你不可认为，不仅是与（创造）世界有关的工程是由圣子完成的，而且自从那时以来，凡是天主所做的一切也是如此。因为\Quote{父爱子，并把一切交在祂手中}，祂确实从起初就爱子，并从最初就把一切交给了祂。因此经上记载：\Quote{在起初，圣言已与天主同在，圣言就是天主；}\BibleRef{若 1:1} 父将\Quote{天上地下的全部权能都交给了祂；}\BibleRef{玛 28:18}\Quote{父不审判任何人，却把一切审判都交给了子}\BibleRef{若 5:22}——甚至从起初就是如此。因为当祂说到全部权能和全部审判，并说万物都是藉着祂造的，万物都已交在祂手里时，祂不容许时间上的任何例外，因为除非它们是\Emph{一切时间}的事物，否则它们就不是\Emph{万物}。因此，是从起初就施行审判、推倒傲慢的高塔、分裂语言、用洪水的暴力惩罚整个世界、向\Place{索多玛}和\Place{哈摩辣}降下火与硫磺的子，祂是来自主的主。因为祂时常下来与人交谈，从\Person{亚当}到族长和先知，藉着异象、梦、镜子、谜语；从起初就为祂的\Emph{经世计划}的进程奠定基础，这计划祂打算贯彻到底。就这样，祂甚至作为天主学习与人在地上交谈，祂就是那将要成为血肉的圣言。但祂这样做（学习或预演），是为我们铺平信德之路，好让我们更轻易相信天主圣子曾降临世界，如果我们知道在过去的时代也发生过类似的事。因为正如这些事件记在圣经中是为了我们\Emph{和我们的教诲}，同样它们也是为了我们而行——（甚至是我们，我说），\Quote{我们这些万世万代的终结已来到的人。}\BibleRef{格前 10:11} 就这样，即使在那时，祂也完全知道人的情感和感受，因为祂总是打算实际取人的组成成分，即身体和灵魂，（例如）询问\Person{亚当}（仿佛不知道）：\Quote{亚当，你在哪里？}\BibleRef{创 3:9}——后悔造了人，仿佛缺乏预见；\BibleRef{创 6:6}试探\Person{亚巴郎}，仿佛不知道人心里想什么；对人发怒，然后又与他们和好；以及其他异端者（在假设中）抓住的、认为不配于天主的一切（软弱和不完美），以贬低创造主，却未考虑这些\Emph{情况}对于圣子来说是足够适合的，因为祂将来要经历人的苦难——饥饿、口渴、眼泪、实际出生和真正死亡，并且为了这样的救恩计划，\Quote{被父贬抑得比天使微小一点}。但你可以肯定，异端者不会承认这些事甚至适合天主圣子，而你们却把它们归于圣父本身，当你们假装祂为了我们使自己变小；而圣经告诉我们，那被贬抑的是由另一位如此对待，而非祂自己。再者，如果祂是\Emph{一位}被\Quote{加冕以光荣和尊荣}，而\Emph{另一位}（即子由父）是加冕者，那又如何？此外，怎么发生这样的事：全能的不可见的天主，\Quote{没有人看见过，也不能看见；祂住在不可接近的光中；}\BibleRef{弟前 6:16}\Quote{祂不住在人手所造的殿宇中；}\BibleRef{宗 17:24}\Quote{在祂面前大地战栗，山岳像蜡一样熔化；}祂把整个世界握在手中\Quote{像鸟巢；}\BibleRef{依 10:14}\Quote{天是祂的宝座，地是祂的脚凳；}\BibleRef{依 66:1} 在祂之内是每一处，而祂自己却不在任何地方；祂是宇宙的极致——我说，祂（虽然是）至高者，怎能在傍晚凉爽时分漫步在乐园中寻觅\Person{亚当}？怎样在\Person{诺厄}进入方舟后关上了门？怎样在\Person{亚巴郎}的帐幕旁、在一棵橡树下休息？怎样从荆棘火焰中呼唤\Person{梅瑟}？怎样在\Place{巴比伦}王的火窑中显现为\Quote{第四位}（虽然那里祂被称为人子）——除非所有这些事件都作为影像、作为镜子、作为谜语（提示未来的降生成人）而发生的呢？确实，即使是这些事，若不是圣经赐给我们，连对天主圣子本身也未必能相信；或许它们即使记在圣经中，也未必能相信是圣父做的事，因为这些异端者把祂带进\Person{玛利亚}的子宫，把祂放在\Person{比拉多}的审判台前，把祂埋在\Person{若瑟}的坟墓里。因此，他们的错误就显明了；因为，他们不知道神圣管理的整个秩序从一开始就是通过圣子的中介而进行的，他们相信圣父本身真被看见、与人交谈、工作、口渴、忍受饥饿（尽管先知说：\Quote{永恒的天主，\Emph{主，地极的创造者}，从不口渴，也不饥饿；}\BibleRef{依 40:28} 更不用说任何时候死亡或被埋葬了！），因此他们相信一直是同一位天主，即圣父，在一切时代亲自做了那些实际上由圣子中介所做的事情。\par

% structure: p0035 source=h2 target=section reason=relative-heading-rank; base=h2

\section{第十七章 描述神性的庄严称号，应用于圣子，而非如\Person{普拉克西亚斯}所主张的仅应用于圣父}

他们更容易认为圣父以圣子的名义行事，而非圣子以圣父的名义行事，尽管主亲自说：\Quote{我因我父的名而来；}\BibleRef{若 5:43}且向圣父说：\Quote{我已将你的名显示给这些人；}\BibleRef{若 17:6}经上同样说：\Quote{奉主名来的，当受赞美；}即圣子以圣父的名而来。至于圣父的名号：全能的天主、至高者、万军之主、以色列的君王、\Quote{自有者}，我们说（因为经上如此教导我们）这些名号也合宜地属于圣子，圣子受这些称呼，并始终在其中行事，从而将它们在他内显示给人。他说：\Quote{凡父所有的，都是我的。}\BibleRef{若 16:15}那么，为何他的名号不也是呢？因此，当你读到全能的天主、至高者、万军之主、以色列的君王、\Quote{自有者}时，要考虑圣子是否也由这些称呼所指明。他凭自身即是全能的天主，因他是全能天主之圣言，并已获得统管万有的权能；他是至高者，因他\Quote{被高举到天主的右边}，如伯多禄在宗徒大事录中所说；\BibleRef{宗 2:22}他是万军之主，因万有都由圣父屈服于他；他是以色列之君王，因那民族的命运特别交托给他；同样他是\Quote{自有者}，因为有许多被称为儿子的，\Emph{却不是}。至于他们所坚持的，即基督之名也属于圣父，他们将在适当的地方听到（我所说的话）。同时，让我立即针对他们从若望默示录中引用的论点作回应：\Quote{我是今在、昔在、将来永在的全能者；}\BibleRef{默 1:8}以及其他所有他们认为使全能天主的名号不适合于圣子的段落。仿佛\Emph{将来永在者}不是全能者；而事实上，全能者的圣子正如天主的儿子是天主一样，同样是全能者。\par

% structure: p0037 source=h2 target=section reason=relative-heading-rank; base=h2

\section{第十八章  先知著作中对唯一天主的称谓。旨在反对异教偶像崇拜，并不排除天主圣子的对应概念。圣子在圣父内}

但阻碍他们轻易察觉圣父名号在圣子内共有这一点的，是圣经每当断定天主仅有一位时的陈述；好像同一圣经没有同时设定两位——既是天主又是主——如我们上面所展示的。他们的论点是：既然我们发现两位和一位，那么两者——圣父与圣子——就是同一且相同的。然而圣经并不需要任何人的论据来辅助，以免显得自相矛盾。它自有其方法，既陈述只有一位天主，又表明有两位——圣父与圣子——并且与自身一致。显然，圣子是被其所提及的。因为，在没有损害圣子的情况下，圣经完全有能力正确地断定天主只有一位，圣子属于这位天主；因为拥有圣子的人并不因此就不再存在——他自己就是唯一的一位，也就是说，凭自身，每当他被提及时不连同圣子。而他在不连同圣子的情况下被提及，是当他被界定为神性的本原，具有\OriginalTerm{第一位}{the first Person}的特征时；这特征必须在圣子之名被提及之前先说，因为首先是圣父被承认，在圣父之后才提到圣子。因此，\Quote{只有一位天主，}圣父，\Quote{除我以外没有别的神。}\BibleRef{依 45:5}当他亲作此宣告时，他并不否认圣子，而是说没有别的神；而圣子与圣父并无分别。实际上，你若仔细查看这类陈述之后的上下文，便会发现它们几乎总是明确针对偶像制造者及崇拜者，为要借天主的唯一性驱逐众多假神；这唯一性却拥有圣子；而既然圣子与圣父不可分割、不可分离，那么即使圣子未被提名，他仍应被视作在圣父内。事实上，如果他明确提名圣子，就等于将他分离，直接说：\Quote{除我以外没有别的神，\Emph{除了我的圣子}。}简言之，他会在将圣子从其他中排除后，实际上使圣子成为另一个。假设太阳说：\Quote{我是太阳，除我以外没有别的，除了我的光芒，}你难道不会注意到这样的陈述多么无用，好像光芒本身不被算在太阳内？因此，他说除他以外没有别的神，是针对外邦人和以色列人的偶像崇拜；甚至也针对我们的异端者，他们用言语制造偶像，如同外邦人用手制造一样；也就是说，他们造出另一位天主和另一位基督。因此，当他见证自己的唯一性时，圣父顾及了圣子的利益，免得基督被认为来自另一位天主，而是来自那位早已说过\Quote{我是天主，除我以外没有别的神}者；他向我们显示他是唯一的天主，但与他的圣子同在，\Quote{他独自铺张诸天。}\BibleRef{依 44:24}\par

% structure: p0039 source=h2 target=section reason=relative-heading-rank; base=h2

\section{第十九章  圣子与圣父在创造万物中的合一。这二者在合作中的合一并不反对天主的真合一。它只反对\Person{普拉克塞亚斯}的等同理论}

但是，他们会急忙曲解祂\Emph{\OriginalTerm{单一性}{singleness}}的这一宣告。祂说：\BibleQuote{我独自}，祂说，\BibleQuote{铺张了苍天}。无疑，这是指相对于其他一切权势的\Emph{独自}；祂以此预先驳斥了异端者的猜想，他们主张世界是由各类天使和权势建造的，他们甚至把造物主本身说成要么是一个天使，要么是一个被差遣来形成外在事物（如世界的组成部分）的下属代理者，但同时祂对\Emph{神圣旨意}一无所知。如果现在正是这个意义下祂独自铺张诸天，那么这些异端者何以如此悖谬地采取他们的立场，以致拒绝承认那智慧（她曾说：\BibleQuote{当祂预备上天时，我在祂面前？} \BibleRef{箴 8:27}）的单一性？——尽管宗徒提问：\BibleQuote{谁知道主的心意？谁做过祂的参谋？} \BibleRef{罗 11:34}，这当然是指除了那与祂同在的智慧以外。\BibleRef{箴 8:30}无论如何，智慧是在祂内、与祂一同构造宇宙，祂并不无知于她所造之物。然而，\Quote{除了智慧}这一短语，与\Quote{除了圣子}（即基督，\BibleQuote{天主的智慧和德能} \BibleRef{格前 1:24}）具有完全相同的含义，正如宗徒所说，祂唯独知道圣父的旨意。\BibleQuote{因为除了那在祂内的圣神外，谁知道在天主内的事？} \BibleRef{格前 2:11}请注意，不是\Emph{没有}祂。所以，有一位使天主不再独处，除了从其他一切\Emph{神祇}处\Quote{独处}。但是（我们若追随异端者），连福音本身也必须被拒绝，因为它告诉我们：万物是由天主透过圣言造成的，没有祂，什么也没有造成。\BibleRef{若 1:3}如果我没有记错，还有另一段经文写道：\Quote{借主的话造了诸天，借祂口中的圣神造了万军。}而今，这圣言，天主的德能和天主的智慧，必是圣子本身。因此，如果祂透过圣子造化万物，那么祂必是透过圣子铺张诸天，因此并非独自铺张，除非就祂是从一切其他神祇处\Quote{独处}（且分别）的意义而言。为此，祂紧接着论及圣子说：\BibleQuote{谁又使说假话者的预兆落空，使占卜者疯狂，使智慧人退后，使他们的知识成为愚拙，并证实祂圣子的言语？} \BibleRef{依 44:25}——例如，当祂说：\BibleQuote{这是我的爱子，我所喜悦的；你们要听祂！} \BibleRef{玛 3:17}祂这样将圣子联于自身，就成为自己解释者，说明祂在什么意义上独自铺张诸天，即意味着\Emph{与祂的圣子一同独处}，因为祂与圣子原为一体。因此，同样这陈述也是圣子的：\BibleQuote{我独自铺张了诸天} \BibleRef{依 44:24}，因为诸天是\Emph{借圣言}建立的。所以，既然当智慧临于圣言时，天被预备，且万物都由圣言造成，那么说连圣子也独自铺张了天，是完全正确的，因为唯独祂为圣父的工作服务。也必须是指祂说：\Quote{我是元始，并到永远，我是自有者。}无疑，圣言是在万物之前。\BibleQuote{在起初已有圣言} \BibleRef{若 1:1}，并且在那起初，祂被圣父差遣。然而，圣父没有起初，因为祂不是由谁生出的；祂也不能被看见，因为祂未受生。那位永远是独一者，不可能有次序或等级。因此，如果他们认定圣父与圣子应被视为同一，其明确目的是为了维护天主的合一性，那么祂的合一性仍完整保留；因为祂是独一的，然而祂有一个圣子，这位圣子与祂同等，同被记载于同一圣经中。既然他们不愿承认圣子是\Emph{\OriginalTerm{独立的位格}{distinct Person}}，次于圣父，唯恐因次而被说成两位神，我们上面已经证明，圣经中实际描述两位为天主和主。为防止他们对此事实反感，我们给出一个理由，说明为什么他们不被说是两位神和两位主，而是两位，即圣父和圣子；这并不是由于他们本体分离，而是由于经世，在此经世中我们宣认圣子与圣父不可分、不可离——在程度上区别，而非在本身上。并且，尽管当分别被称呼时，祂被称为天主，祂并不因此造成两位神，而是一位；这正是由于祂从与圣父的合一中，有资格被称为天主。\par

% structure: p0041 source=h2 target=section reason=relative-heading-rank; base=h2

\section{第二十章 普拉克塞阿斯用来支持其异端的圣经章节为数甚少。德尔图良提及它们。}

但是，我必须更加费心来驳斥他们的论点，因为他们从圣经中挑选一些章节来支持自己的意见，而拒绝考虑其他显然维护信仰准则且丝毫不破坏天主性合一、并完全承认\OriginalTerm{单一君主制}{Monarchy}的要点。因为在旧约圣经中，他们抓住无他，唯有：\BibleQuote{我是天主，在我以外没有别神；} \BibleRef{依 45:5}同样，在福音中，他们只注意主对斐理伯的回答：\BibleQuote{我与圣父原是一体} \BibleRef{若 10:30}，以及\BibleQuote{看见了我，就是看见了圣父；我在圣父内，圣父在我内} \BibleRef{若 14:9}。他们要使两约的全部启示都顺服于这三段经文，然而唯一正确的途径是根据多数章节来理解少数章节。但在争辩中，他们不过是按所有异端者的原则行事。因为，既然在总体的大量经文里只找到少数（对他们有利的）见证，他们便固执地用少数对抗多数，并以后来的（解释）对抗先前的。然而，那从起初为每一情况所制定的准则，就以其定则反对后来的\Emph{揣测}，正如它也反对较少的（见证）一样。\par

% structure: p0043 source=h2 target=section reason=relative-heading-rank; base=h2

\section{第21章 在本章及随后四章中，通过详细分析\BibleBook{若望}{福音}表明，圣父与圣子常被说成\OriginalTerm{不同位格}{Distinct Persons}}

因此，请想一想，在这位斐理伯提出问题之前，在你进行任何讨论之前，在这同一部福音中有多少段落向你提出了它们的规范权威。首先，若望为他福音所写的序言立即呈现在手，它向我们显示了那将要成为血肉者先前是什么。\BibleQuote{在起初已有圣言，圣言与天主同在，圣言就是天主。他在起初与天主同在：万物是藉着他而造成的；没有他，什么也没有造成。}\BibleRef{若 1:1}既然这些话只能按其所写的来理解，那么毫无疑问显示了有一位从起初就有的，还有一位祂始终与他同在：一个是天主的圣言，另一个是天主（尽管圣言也是天主，但这里的\Emph{天主}被视为天主子，而非圣父）；一位是万物藉着他而有的，另一位是万物由他而有的。但我们在什么意义上称祂为另一位，我们已经多次描述过。在我们称祂为另一位时，必须意味着祂并非同一——并非同一，却也不是分离的；是藉着\OriginalTerm{安排}{dispensation}而不同的，不是藉着分裂。因此，那成为血肉者与圣言所由来的那一位并非完全相同。\BibleQuote{我们看见了祂的光荣——正如父独生者的光荣，}\BibleRef{若 1:14}（请注意）不是父的光荣。祂\Quote{宣明了}（那独一）\Quote{父怀中的}事；父并没有泄露自己怀中的秘密。因为在这之前有另一句话：\Quote{从来没有人见过天主。}然后，当祂被若翰（洗礼者）称为\Quote{天主的羔羊}时，\BibleRef{若 1:29}祂并非被描述为与祂所是的那位（即祂是爱子）相同。无疑，祂永远是\Emph{天主子}，但祂并非祂所从出的那一位本身。纳塔乃耳立刻在祂身上认出了这一（神圣关系），\BibleRef{若 1:49}正如伯多禄在另一场合所做的：\Quote{你是天主子。}\BibleRef{玛 16:16}祂也亲自肯定他们的信念是正确的，因为祂回答纳塔乃耳说：\Quote{因为我对你说，我看见你在无花果树下，你就信了吗？}\BibleRef{若 1:50}同样，祂称伯多禄为\Quote{有福的}，因为\Quote{不是血肉启示了他}——即他领悟了圣父——\Quote{而是我在天之父。}\BibleRef{玛 16:17}藉着断言这一切，祂确定了两个位格之间的区别：即当时在地上的子——伯多禄已宣认为天主子；以及在天上的父——祂启示给伯多禄他所发现的，即基督是天主子。当祂进入圣殿时，祂称它为\Quote{我父的殿宇，}\BibleRef{若 2:16}是作为子而说的。在祂对尼苛德摩的讲话中，祂说：\BibleQuote{天主竟这样爱了世界，甚至赐下了自己的独生子，使凡信祂的人不至丧亡，反而获得永生。}\BibleRef{若 3:16}又说：\BibleQuote{因为天主没有派遣子到世界上来审判世界，而是为叫世界藉着他而获救。那信从他的，不受审判；那不信的，已受了审判，因为他没有信从天主独生子的名字。}\BibleRef{若 3:17}此外，当若翰（洗礼者）被问及他对耶稣所知为何时，他说：\BibleQuote{父爱子，并将一切交在他手中。那信从子的，便有永生；那不信从子的，不会见到生命，天主的义怒常在他身上。}\BibleRef{若 3:35}实际上，祂向撒玛黎雅妇人启示了谁？不就是那称为基督的默西亚吗？\BibleRef{若 4:25}因此祂当然表明了自己不是父，而是子；在别处祂也被明确称为\Quote{基督、天主子，}\BibleRef{若 20:31}而非父。因此祂说：\BibleQuote{我的食物就是承行派遣我者的旨意，并完成祂的工程。}\BibleRef{若 4:34}而对于犹太人，关于治好软弱者的事，祂说：\Quote{我父一直工作到现在，我也工作。}\BibleRef{若 5:17}\Quote{我父和我}——这是子的话。正是为此，\Quote{犹太人越发想要杀害祂，不但因为祂破坏了安息日，而且因为祂称天主是自己的父，使自己与天主平等。}于是祂回答他们说：\BibleQuote{子不能自己做什么，只有看见父做什么，才能做什么；凡父所做的，子也照样做。因为父爱子，凡自己所做的都指示给祂；并且还要把比这些更大的工作指示给祂，叫你们惊奇。就如\Emph{父}唤回死者并使他们复活，照样子也使他所愿意的人复活。父不审判任何人，但把所有审判权全交给了子，为叫众人尊敬子如同尊敬父。那不尊敬子的，就是不尊敬派遣祂来的父。我实实在在告诉你们：那听我话，并信派遣我来者的，便有永生，不受审判，而已出死入生。我实实在在告诉你们：时候要到，且现在就是，死者要听见天主子的声音，凡听见的，就必生存。就如父在自己内有生命，照样也赐给子在自己内有生命；并且赐给祂行审判的权柄，因为祂是人子，}\BibleRef{若 5:19}这是按肉体而言，正如祂也藉着祂的圣神而成为天主子。随后祂继续说：\Quote{但我有比若翰更大的证据；因为父交给我要完成的工程，就是我所行的这些工程，为我作证：是父派遣了我。派遣我来的父，也亲自为我作证。}\BibleRef{若 5:36}但祂立刻补充说：\Quote{你们从来没有听过祂的声音，也从来没有见过祂的形状。}从而肯定了在以前的时代，并不是父，而是子常被看见和听到。最后祂说：\Quote{我因我父的名而来，你们却不接纳我。}因此，我们读到的那一位总是子，以全能、至高天主、君王和主的称号出现。对于那些问\Quote{我们该做什么，才算作天主的事业}的人，\BibleRef{若 6:29}祂回答说：\Quote{\Emph{天主的事业}，就是要你们信从祂所派遣来的那一位。}祂也宣称自己是\Quote{从天降下的父所赐的食粮}；并补充说，\Quote{凡父赐给祂的，必到祂这里来，而祂必不拒绝他们，因为祂从天降下，不是为承行自己的旨意，而是为承行那派遣祂者的旨意；父的旨意就是：凡看见子，并信从子的，必获得永生，并且在末日复活。没有人能到祂这里来，除非父吸引他；凡是听父的教导而学习的，都到祂这里来。}\BibleRef{若 6:37}然后祂明确地说：\Quote{这不是说有人见过父}；从而向我们表明，人是藉着父的圣言而受训诲和教导的。然后，当许多人离开祂，祂转向宗徒们问他们是否也要走时，西满伯多禄的回答是什么？\Quote{主！唯你有永生的话，我们相信而且知道你是天主的圣者。}（现在告诉我，他们相信）祂是父，还是父的基督？\par

% structure: p0045 source=h2 target=section reason=relative-heading-rank; base=h2

\section{第二十二章 所引\Person{圣若望}的各项段落，以显示圣父与圣子之间的区别。甚至\Person{普拉克塞斯}的经典文本——\BibleQuote{我与父原是一体}——也显示为反对他。}

再者，祂所宣布的是谁的教训，叫众人都惊奇？是祂自己的，还是父的？因此，当他们彼此疑惑祂是否基督（当然不是作为父，而是作为子）时，祂对他们说：\Quote{你们并非不知道我从哪里来；我并不是凭自己来的，但那派我来的乃是真实的，你们不认识祂；我却认识祂，因为我是从祂来的。}祂并没有说：因为我本身就是祂；我自己派了我自己；而是说：\Quote{祂派了我来。}同样，当法利塞人派人去捉拿祂时，祂说：\Quote{我还有不多的时候与你们同在，然后我往派我来的那一位那里去。}然而，当祂宣称祂不是独自一人，并说了这些话：\Quote{我与派我来的父}，\BibleRef{若 8:16}难道祂不是表明有两位——两位，且不可分吗？事实上，这正是祂教训他们的总纲和实质，即他们是不可分离的两位；因为，祂援引律法——它肯定两个人的见证是真实的——之后，立刻加上：\Quote{我是为自己作证的一位；还有派我来的父也为我作证。}现在，如果祂是一——同时是子和父——祂断不会援引律法的认可，它要求的不是一个人的见证，而是两个人的。同样，当他们问祂父在哪里时，祂回答他们说，他们既不认识祂，也不认识父；在这答复中，祂清楚告诉他们有\Emph{两位}，他们并不认识。即使说\Quote{他们若认识祂，也就认识父}，这当然不意味着祂自己既是父又是子；而是因为这两位不可分离，若没有另一位，其中一位不可能被认识或不被认识。祂说：\Quote{派我来的那位是真实的；我把从祂那里听到的告诉世人。}\BibleRef{若 8:26}圣经记述接着以浅显的方式解释说\Quote{他们不明白祂是对他们讲论父}，虽然他们本应知道父的话是在子内\Emph{说出来}的，因为他们读到耶肋米亚书说：\Quote{上主对我说：看，我将我的话放在你口中；}\BibleRef{耶 1:9}以及依撒意亚书说：\Quote{上主赐给了我受教的口舌，为使我晓得应怎样及时发言。}\BibleRef{依 50:4}据此，\Emph{基督}亲自说：\Quote{那时你们就会知道我即是祂，并且我什么也不凭自己说；而是父怎样教导了我，我就怎样讲，因为派我来的那位与我同在。}\BibleRef{若 8:28}这同样证明他们是两位，（虽）不可分。同样，当祂在辩论中责备犹太人，因为他们想要杀祂时，祂说：\Quote{我讲的是我在父那里所看见的，你们做的却是从你们的父那里所看见的；}\Quote{但如今你们想要杀我，一个将我从天主那里听到的真理告诉你们的人；}又说：\Quote{如果天主是你们的父，你们就会爱我，因为我出自天主并由天主而来，}（然而他们并不因此被分离，尽管祂宣称祂是\Emph{从父}出来的。有些人确实抓住这些话提供的机会\Emph{来提出祂分离的异端}；但祂从天主出来，好比光线从太阳发出，河流从泉源流出，树木从种子长出）；\Quote{我没有附魔，我尊敬我的父；}又说：\Quote{如果我光荣自己，我的光荣算不了什么；光荣我的是我的父，就是你们所说的你们的天主。你们不认识祂，我却认识祂；我若说不认识祂，我就像你们一样是个说谎者；但我认识祂，也遵守祂的话。}\BibleRef{若 8:54}但当祂接着说：\Quote{你们的\Emph{祖先}亚巴郎欢欣地看到我的日子，他看见了，且喜乐。}祂确实证明不是父显现给亚巴郎，而是子。同样，祂关于那个\Emph{生来}瞎眼的人宣告说：\Quote{祂必须做那派祂来的父的工；}\BibleRef{若 9:4}并且祂使那人看见后，对他说：\Quote{你信天主子吗？}然后，当那人问祂是谁时，祂就向他启示自己，就是祂曾向他宣布为信德正当对象的那个\Emph{天主的}子。在后一段经文\BibleRef{若 10:15}祂宣称祂被父认识，父也被祂认识；并说祂如此全然为父所爱，以至于祂舍去自己的性命，因为祂从父领受了这一命令。当犹太人问祂是否是那位基督（当然是指天主的基督；因为直到今日，犹太人期待的并非父本身，而是天主的基督，从未说父会作为基督来临），祂对他们说：\Quote{我告诉你们，你们却不信：我以我父的名所做的工，这些工为我作证。}见证什么？正是他们所询问的事——祂是否是那位天主的基督。接着，又关于祂的羊，以及（保证）没有人能把它们从祂手里夺去，祂说：\Quote{我的父，把羊赐给我的那位，比一切更大；}立刻又说：\Quote{我与父是一体的。}在此，他们就站稳了立场，过于执迷，不，过于盲目，看不到首先这经文暗示了两位的存在——\Quote{\Emph{我与父}}；其次，有一个复数谓语\Quote{\Emph{是}}，不能只用于一位；最后，（谓语终止于一个抽象名词，而非人称名词）——\Quote{我们是一\Emph{体}}\OriginalTerm{一}{Unum}，而非\Quote{一位}\OriginalTerm{一个}{Unus}。因为如果祂说了\Quote{一位}，或许会对他们的意见有所支持。\Emph{Unus}无疑表示单数；但（在此情况下）\Quote{两位}仍然是阳性主语。因此祂说\Emph{Unum}，一个中性词，不表示数字的单一，而是表示本质、相似、结合、父对子的爱——父爱子，以及子对父旨意的服从。当祂说：\Quote{我与父是一体}\Emph{本质上}——\OriginalTerm{一}{Unum}——祂表明有两位，祂使他们平等并合而为一。因此祂紧接着这句话，说祂\Quote{从父那里向他们显示了许多善工}，没有一件工祂该受石击。\BibleRef{若 10:32}为了防止他们认为祂该受此惩罚，好像祂曾因说\Quote{我与父是一体}而自认为就是天主本身，即父，祂表明自己是父的神圣之子，而非天主本身，祂说：\Quote{如果你们的法律上写着：我说过你们是神；并且圣经不能废弃，那么父所祝圣并派到世界来的那位，你们怎可以说祂亵渎，因为祂说：我是天主子？我若不行我父的工，你们就不必信我；但我若行了，你们纵然不信我，也该信这些工，以便知道父在我内，我在父内。}因此，必然是藉着这些工，父在子内，子在父内；同样，藉着这些工，我们明白父与子\Emph{是一}。因此祂始终努力达到这一结论：尽管他们拥有同一能力和本质，他们仍应被相信是两位；否则，除非他们被相信是两位，子就不可能被相信有任何存在。\par

% structure: p0047 source=h2 target=section reason=relative-heading-rank; base=h2

\section{第二十三章 同一福音中更多经文以证明同一部分公教信仰。\Person{普拉克瑟斯}关于崇拜两位天主的嘲讽被驳斥。}

又，当\Person{玛尔大}在后来的段落中承认祂是天主子时，\BibleRef{若 11:27}她并不比\Person{伯多禄}\BibleRef{玛 16:16}和\Person{纳塔乃耳}\BibleRef{若 1:49}更会犯错；然而，即使她犯了错，她也立刻会学到真理：因为，看，当主要从死者中复活她兄弟时，主举目向天，对圣父说话——当然是作为圣子说：\BibleQuote{父啊，我感谢祢，因为祢总是俯听我；正是为了这些站在旁边的人，我才对祢说了这话，好叫他们相信祢派遣了我。}\BibleRef{若 11:41}但（后来）在祂灵魂烦乱时，祂说：\BibleQuote{我该说什么呢？父啊，救我脱离这时辰；但正是为此我才到了这时辰；父啊，光荣祢的名吧！}\BibleRef{若 12:27}——祂是以圣子的身份说的。（另有一次）祂说：\BibleQuote{我因我圣父的名而来。}\BibleRef{若 5:43}因此，圣子的声音（向圣父）本是独自足够的。但是，看，圣父从天上以大量（证据）回答，为给圣子作证：\BibleQuote{这是我的爱子，我所喜悦的，你们要听从祂。}\BibleRef{玛 17:5}同样，在那\Emph{断言}中：\BibleQuote{我已光荣了，还要再光荣}\BibleRef{若 12:28}，你发现有多少位格，顽固的\Person{普拉克瑟斯}？难道不是如同声音的数量那样多吗？你有圣子在地上，你有圣父在天上。这并不是分离；这不过是神圣的时期。然而我们知道，天主在无底深处，且无处不在；但那是在能力和权柄上。我们也确信，圣子与祂不可分，处处与祂同在。然而，在\OriginalTerm{经济}{Economy}或\OriginalTerm{时期}{Dispensation}本身中，圣父愿意圣子被视为在地上，而祂自己在天上；圣子自己也举目向天，祈祷，并向圣父恳求；祂也教导我们举起自己，祈祷：\BibleQuote{我们在天的父}\BibleRef{玛 6:9}等等——尽管祂确是无所不在的。这个\Emph{天}，圣父愿意作为祂自己的宝座；而祂使圣子\Quote{稍微逊于天使}，将祂派遣到地上，但同时也要\Quote{以光荣和尊贵给祂加冕}，甚至通过将祂带回天上。祂现在使这应验，当祂说：\BibleQuote{我已光荣了\Emph{祢}，还要再光荣\Emph{祢}。}圣子从地上提出请求，圣父从天上给予应许。那么，你为何使圣父和圣子都成为说谎者？如果圣父从天上对圣子说话，而祂自己就是地上的圣子，或者圣子向圣父祈祷，而祂自己就是天上的圣子，那么圣子如何向圣父请求祂自己，既然圣子就是圣父？或者，另一方面，圣父如何向自己应许，既然圣父就是圣子？即使我们主张他们是两位不同的神，正如你喜欢向我们指控的那样，也比维护你这样多变善变的天主更可容忍！因此，主在面前的经文中向在场的人宣告：\BibleQuote{这声音不是为我而来，而是为你们}\BibleRef{若 12:30}，为叫他们也同样相信圣父和圣子，各自在自己的名字、位格和位置中。\BibleQuote{耶稣又呼喊说：信我的，不是信我，而是信那派遣我的}\BibleRef{若 12:44}，因为人藉着圣子相信圣父，而圣父也是信仰出于圣子的权威。\BibleQuote{看见我的，就是看见那派遣我的。}如何如此？正因为（如祂后来宣告的）：\BibleQuote{我没有凭自己说话，而是派遣我的父：祂给了我命令，我该说什么，该讲什么。}\BibleRef{若 12:49}因为\BibleQuote{上主天主赐给了我受教的舌头，使我知道我该何时说话}\BibleRef{依 50:4}——我所实际说的言语。\BibleQuote{正如父对我所说的，我就这样讲。}\BibleRef{若 12:50}现在，这些事是如何对祂说的，福音作者和爱徒\Person{若望}比\Person{普拉克瑟斯}更清楚；因此他按自己的意思补充说：\BibleQuote{在逾越节庆日前，耶稣知道圣父已将一切交在祂手中，并且祂是从天主而来，要往天主那里去。}然而，\Person{普拉克瑟斯}却认为，是圣父从祂自己发出，又回到祂自己；因此魔鬼放在\Person{犹达斯}心里的背叛，不是对圣子的背叛，而是对圣父本身的背叛。但就此而言，事情对魔鬼或异端者都无好结果；因为，即使在圣子的事上，魔鬼对祂所行的背叛也未给他带来任何益处。因此，是那位在人子内的天主子被出卖了，正如圣经后来所说：\BibleQuote{如今人子受到了光荣，天主也在祂身上受到了光荣。}这里\Quote{天主}指的是谁？一定不是圣父，而是圣父的圣言，祂在人子内——即在肉身内，\Person{耶稣}已藉着\Emph{神圣的}能力和圣言得了光荣。\BibleQuote{而且天主，}祂说，\BibleQuote{也要在祂自己内光荣祂；}也就是说，圣父要光荣圣子，因为祂在自己内拥有祂；即使祂被屈辱于地，置于死地，祂也将很快藉着复活光荣祂，并使祂成为战胜死亡者。\par

% structure: p0049 source=h2 target=section reason=relative-heading-rank; base=h2

\section{第二十四章 论\Person{圣斐理伯}与基督的对话。看见了我，就是看见了父。以反\Person{普拉克瑟斯}的意义解释此经文。}

但即使在那时，仍有一些人不明白。因为多默长久不信，便说：\BibleQuote{主，我们不知道祢往哪里去，又怎能知道那条路呢？耶稣对他说：我就是道路、真理、生命，除非经过我，谁也不能到圣父那里去。如果你们认识我，也就认识我的父；从现在起，你们认识祂，并且已经看见祂了。} \BibleRef{若 14:5} 现在我们谈到斐理伯，他因期望看见圣父而激动，却不明白如何理解\Quote{看见圣父}，便说：\Quote{请把圣父显示给我们，我们就心满意足了。} 主就回答他说：\Quote{斐理伯，这么长久我与你们在一起，你还不认识我吗？} 祂说他们应该认识的是谁？——这正是讨论的唯一焦点。他们应该认识祂是圣父呢，还是圣子？如果祂是圣父，普拉克色亚就必须告诉我们，基督既已长久与他们在一起，怎么可能曾被（我不说理解，甚至）认为是圣父。在所有圣经中，祂被清楚地向我们指明——旧约中是天主的基督，新约中是天主的圣子。正是以这种身份，祂在古时被预言；也正是以这种身份，连基督自己都宣告祂；不，连圣父自己也是如此，祂从天上公开承认祂为自己的圣子，并作为圣子光荣祂。\BibleQuote{这是我的爱子；} \BibleQuote{我已经光荣了祂，我还要光荣祂。} 也是以这种身份，祂被门徒们信从，被犹太人拒绝。而且，正是以这种身份，每当祂提到圣父、尊崇圣父、光荣圣父时，祂希望他们接纳祂。既然如此，祂与他们长久交往后，他们所不认识的不是圣父，而是圣子。因此，主责备\Emph{斐理伯}不认识祂——这正是他们无知的对象——祂愿意自己被承认为那个祂责备他们长久以来所不认识的那一位，简言之，就是圣子。现在可以看出，\BibleQuote{看见我的，就是看见了圣父} \BibleRef{若 14:9} 这句话在什么意义上说的——正是在前一段\BibleQuote{我与父原是一体} \BibleRef{若 10:30} 的同一意义上。为什么？因为\BibleQuote{我从父出来，到了\Emph{世界上}} \BibleRef{若 16:28}，\BibleQuote{我是道路，除非经过我，谁也不能到父那里去} \BibleRef{若 14:6}，\BibleQuote{除非父吸引他，没有人能到我这里来} \BibleRef{若 6:44}，\BibleQuote{一切由父交给了我} \BibleRef{玛 11:27}，\BibleQuote{父怎样复活（死者），子也照样复活} \BibleRef{若 5:21}，以及\BibleQuote{如果你们认识我，也就认识我的父} \BibleRef{若 14:7}。因为在所有这些段落中，祂表明自己是圣父的使者，藉着祂的中介，圣父可以在祂的工程中被看见，在祂的话语中被听见，并在圣子管理圣父言行的事工中被认识。圣父确实不可见，正如斐理伯在法律中所学到的，他当时应当记得：\BibleQuote{没有人看见天主还能生活} \BibleRef{出 33:20}。所以他因渴望看见圣父，仿佛祂是可见的存在而受责备，并被教导说：圣父只在圣子内，藉祂大能的作为而成为可见的，并非在祂位格的显现中。如果祂确实要人把圣父理解为与圣子相同，便说\BibleQuote{看见我的，就是看见了圣父}，那么为何祂紧接着又说：\BibleQuote{你不相信我在父内，父在我内吗？} \BibleRef{若 14:10} 祂本应该说：\Quote{你不相信我就是父吗？} 祂如此强调这一点，若不是为了澄清祂希望人明白的事——即祂是圣子——又有什么其他目的呢？接着，祂又说：\BibleQuote{你们不相信我在父内，父在我内吗？} \BibleRef{若 14:11} 祂之所以加重这个问题的语气，正是为了不让人们因祂曾说\BibleQuote{看见我的，就是看见了圣父}而以为祂是圣父；因为祂从不希望自己被如此看待，一直声称自己是圣子，并从圣父而来。然后，祂还最清楚地指明了两位的合一，为使人不再想看见圣父仿佛祂可被单独看见，并使圣子被视为圣父的代表。然而，祂并没有忽略解释圣父如何在圣子内、圣子在圣父内。祂说：\BibleQuote{我对你们所说的话，不是凭我自己说的} \BibleRef{若 14:10}，因为这些确实是圣父的话；\BibleQuote{是住在我内的父在做祂的工作。} 因此，正是藉着祂大能的作为和祂教义的话语，住在圣子内的圣父使自己可见——甚至藉着那些祂住在其中、也藉着祂所住的那一位的\Emph{言与行}；两位的专属特性因祂说\Quote{我在父内，父在我内}这件事而显明。因此祂又加上：\Quote{你们应相信——} 什么？相信我是父？我找不到这样写，而是\Quote{相信我在父\Emph{内}，父在我内；否则，因着那些工作而相信我。} 意思是指那些工作，藉此圣父显明自己住在圣子内，不是为人的眼光，而是为人的理智。\par

% structure: p0051 source=h2 target=section reason=relative-heading-rank; base=h2

\section{第25章　护慰者，或圣神。祂与圣父和圣子位格上不同，但天主性上合一不可分。出自圣若望福音的其他引文。}

继\Person{斐理伯}的询问之后，以及\Emph{主}对它的全部处理，直到\Emph{若望}福音的末尾，继续向我们提供同类陈述，区分圣父和圣子，并各具属性。此外，还有护慰者\Emph{或安慰者}，祂应许向圣父祈求，并在祂升到圣父那里之后从天上派遣。\Emph{祂被称为} \Quote{另一种护慰者，}确实；见：\BibleRef{若 14:16}但祂如何是\Emph{另一种}，我们已表明；\Quote{祂要从我领受，}基督说，见：\BibleRef{若 16:14}正如\Emph{基督}自己从圣父领受。因此，圣父在圣子内、圣子在护慰者内的联系，产生三位相连的位格，\Emph{他们彼此有别}。这三位是\Emph{\OriginalTerm{本质}{essence}}，而非\Emph{\OriginalTerm{位格}{Person}}，正如所说：\Quote{我与父原是一体，}见：\BibleRef{若 10:30}这是指本体的合一，而非数目的单一。通读整部\Emph{福音}，你会发现，你相信是圣父的那一位（被描述为替圣父行事，尽管你，就你而言，确实以为\Quote{父是园丁，}见：\BibleRef{若 15:1}必定曾在地上）被圣子再次确认为在天上，当祂\Quote{举目向天，}见：\BibleRef{若 17:1}将祂的门徒交托于圣父的保护时，见：\BibleRef{若 17:11}。此外，在另一部福音中，我们有清楚的启示，即\Emph{圣子与圣父的区别}：\Quote{我的天主，我的天主，你为什么舍弃了我？}见：\BibleRef{玛 27:46}以及（在第三部福音中）：\Quote{父啊，我把我的灵魂交在你手中。}见：\BibleRef{路 23:46}但即便没有这些段落，我们在祂复活并荣耀战胜死亡之后，也遇到满意的证据。如今，祂自谦的一切束缚都已除去，祂本可能，如果可能的话，向如此忠信的妇女（如\Person{玛利亚玛达肋纳}）显示自己为圣父，当她出于爱而接近触摸祂时，并非出于好奇，也非带着\Person{多默}的怀疑。\Emph{但并非如此}；\Person{耶稣}对她说：\Quote{你别摸我，因为我还没有升到我的父那里去；你到我的弟兄们那里去}（即使在此，祂证明自己是圣子；因为若祂是圣父，祂会称他们为祂的\Emph{儿女}，而非祂的\Emph{弟兄}），\Quote{对他们说：我升到我的父和你们的父那里去，升到我的天主和你们的天主。}见：\BibleRef{若 20:17}\Emph{那么，这是否意味着，我升到}作为圣父升到圣父，作为天主升到天主？或作为圣子升到圣父，作为圣言升到天主？为什么这部福音在其结尾处暗示这些事曾被写下，难道不是用其自己的话说：\Quote{为叫你们相信耶稣基督是天主子？}见：\BibleRef{若 20:31}因此，每当你取这部福音的任何陈述，并用它们来证明圣父与圣子的同一，以为它们有助于你的观点，你就是在反对这部福音的明确目的。因为这些事当然不是为叫你们相信耶稣基督是圣父，而是圣子。\par

% structure: p0053 source=h2 target=section reason=relative-heading-rank; base=h2

\section{第二十六章 对\BibleBook{玛窦}{福音}和\BibleBook{路加}{福音}的简要引用。它们与\BibleBook{若望}{福音}一致，关于圣父和圣子位格的区别。}

除了斐理伯的谈话和主对它的回答之外，读者将会注意到，我们通览了\BibleBook{若望}{福音}，\Emph{以表明}许多其他含义清晰的经文，无论是在那一章之前还是之后，都严格符合那单一而突出的陈述，并且必须按照所有其他经文（而非与它们对立）以及其本身内在和自然的含义来解释。我在此不会大量引用其他福音书的支持，它们通过主的\Emph{\OriginalTerm{诞生}{nativity}}坚定了我们的信仰；只需指出：那必须由贞女诞生的那位，被天使亲自明确宣告为天主子：\BibleQuote{圣神要临于你，至高者的能力要庇荫你，因此，那要诞生的圣者，将称为\Emph{天主子}。}\BibleRef{路 1:35} 对于这段经文，他们甚至也想吹毛求疵；但真理必将得胜。当然，他们说，天主子就是天主，至高者的能力就是至高者本身。他们毫不迟疑地暗示，如果这是真的，经文就会明写。他如此害怕谁，以至于不直白地宣告：\Quote{天主要临于你，至高者要庇荫你？}现在，通过说\Quote{天主的圣神}（尽管天主的圣神\Emph{就是天主}）而不直接提天主之名，他愿人理解整个\Emph{\OriginalTerm{天主性}{Godhead}}中将要退入\Quote{子}这个称号的那一部分。此处的天主的圣神必定与圣言相同。因为，正如若望说\BibleQuote{圣言成了血肉}\BibleRef{若 1:14}，我们也在提及圣言时理解圣神；同样，在这里，我们也在圣神的名号中承认圣言。因为圣神是圣言的实体，圣言是圣神的运作，二者是\Emph{一}（且同一）。那么，当若望说祂\Quote{成了血肉}时，他必指\Emph{同一位}；而当天使宣告祂\Quote{将要诞生}时，他必指\Emph{另一位}，如果圣神不是圣言，圣言也不是圣神。因为，正如天主的圣言实际上并非那拥有圣言者，同样，圣神（尽管被称为天主）实际上也并非那被称为祂的\Emph{圣神}的那位。属于他者的任何事物，实际上并不等同于他所属于的那位本身。显然，当某物从一个位格主体发出并因此属于他时，由于它从他而来，它在性质上可能恰恰与那发出它并拥有它的位格主体相同。因此，圣神是天主，圣言是天主，因为都出于天主，但并非实际上等同于祂所出的那位。然而，那出自天主的，虽是实际存在者，却不能单独就是天主本身，而只是在同一实体且作为实际存在者和全体的部分的意义上，才是天主。更何况\Quote{至高者的能力}不会是至高者本身，因为它不是实际存在者，作为圣神——就像\Emph{智慧}（天主之）和\Emph{眷顾}（天主之）不是天主一样；这些\Emph{属性}不是实体，而是特定实体的依附体。能力是圣神的附属，但本身不能是圣神。因此，这些事物，无论它们是什么——（我指）天主的圣神、圣言和能力——被赋予贞女后，由她所生的就是天主子。在其他福音书中，祂自己从童年起就见证这就是祂：\BibleQuote{你们不知道我必须在我\Emph{父亲}那里吗？}\BibleRef{路 2:49} 撒殚在诱惑中也知道祂是：\BibleQuote{你如果是\Emph{天主子}。}相应地，魔鬼也承认祂是：\BibleQuote{我们知道你是谁，你是天主的\Emph{圣者}。}祂自己敬拜祂的\Quote{\Emph{父亲}}。当伯多禄承认祂是天主的\BibleQuote{基督（之子）}时，\BibleRef{玛 16:17}祂不否认\Emph{这关系}。祂在圣神内欢跃，向父说：\BibleQuote{父啊，我感谢\Emph{你}，因为你将这些事瞒住了智慧及明达的人。}\BibleRef{玛 11:25} 此外，祂还肯定除了\Emph{子}外，没有人认识父；并应许，作为\Emph{父的子}，祂要在父面前承认那些承认祂的人，并否认那些否认祂的人。\BibleRef{玛 10:32} 祂也引入一个比喻，派遣\Emph{子}（而非父）到葡萄园去，祂是在许多仆人之后被派遣的，\BibleRef{玛 21:33}并被园户杀害，由父报复。祂也不知道最后的日子和时辰，只有父知道。\BibleRef{玛 24:36} 祂将王国赐给门徒，正如祂说，这王国是父为祂预备的。\BibleRef{路 22:29} 祂有权柄，若愿意，可以求父派遣十二军以上的天使来帮助祂。\BibleRef{玛 26:53} 祂呼喊说天主舍弃了祂。\BibleRef{玛 27:46} 祂将灵魂交托在父手里。\BibleRef{路 23:46} 复活后，祂应许门徒必将父的恩许派遣给他们；\BibleRef{路 24:49} 最后，祂命令他们因父及子及圣神之名授洗，而非因一位独一的天主。而且，我们不是一次，而是三次浸入三个位格，每次分别呼号祂们的名。\par

% structure: p0055 source=h2 target=section reason=relative-heading-rank; base=h2

\section{第二十七章 如此确立了圣父与圣子的区别后，他现在证明二性的区别，这二性毫无混淆地结合在圣子位格中；普拉克西亚斯的诡辩因此被揭露。}

但我为何还要在这些如此明显的事上停留，既然我应该攻击他们试图掩盖最清晰证明的那些要点？因为他们在圣父与圣子之间的区分上已被全面驳倒，我们坚持这一区分而不破坏其不可分的结合——正如（借着例子）太阳与光线、泉源与河流——然而借着（他们的妄念）一个不可分的数字，（有着）两种和三种的结果，他们仍试图解释这个\Emph{区分}，以使之附和他们的观点：因此，全都在一个位格内，他们区分出两个：圣父和圣子，理解圣子为肉身，即人，即耶稣；圣父为灵，即天主，即基督。因此，他们虽然主张圣父和圣子是同一的，实际上却始于分裂它们，而非统一它们。因为如果耶稣是一个，基督是另一个，那么圣子将不同于圣父，因为圣子是耶稣，而圣父是基督。我猜想，他们是在\Person{瓦伦提诺}的学校里学到这样的\Emph{\OriginalTerm{独一统治}{monarchy}}，造出两个——耶稣和基督。但事实上，他们的这一观念在我们先前的主张中已被驳倒，因为天主的圣言或天主的圣神也被称为至高者的能力，他们却把这能力认作圣父；然而这些关系并不等同于它们所属的那一位，而是从祂发出并属于祂。然而，在这个异端观点上，另一个驳斥正等着他们。看，他们说，天使曾宣告：\Quote{因此，那由你而生的圣者，将称为天主子。}\BibleRef{路 1:35} 因此，（他们争辩说）既然是肉身诞生了，那么肉身必定是天主子。不，（我回答说）这是指着圣神说的。因为童贞女确实是因圣神怀孕的；祂所孕育的，她生了出来。因此，那被孕育且将要生出的，就必须诞生；也就是说，那圣神，祂的\Quote{名字要称为厄玛奴耳，翻译出来就是：天主与我们同在。}\BibleRef{玛 1:23} 再者，肉身不是天主，所以不能指着肉身说：\Quote{那圣者将称为天主子，}而只能指着那在肉身中出生的神性存在，关于祂，圣咏也说：\Quote{因为天主在其中成为人，并因圣父的旨意建立了它。} 那么是什么神性位格在其中诞生？是圣言，以及因圣父的旨意而与圣言一同降生成人的圣神。因此，圣言成了肉身；这必须是我们探究的要点：圣言如何成了肉身——是通过在肉身中转变形状，还是通过真正穿上肉身。当然是真正地以肉身披戴自己。此外，我们必须相信天主是不变的，无形无相的，因为祂是永恒的。但变形是对先前存在之物的毁灭。因为凡变形为另一物的，就停止是它原先所是的，而开始是它先前所不是的。然而，天主既不停止是祂所是的，也不能成为与祂之所是相异的任何事物。圣言就是天主，而且\Quote{上主的话永远常存}——甚至因在祂自己固有的形态中不变地持守。现在，如果祂不允许被变形，那么祂成为肉身必须被理解为：祂进入在肉身内，借着肉身被显示、被看见、被触摸；因为其他各点同样需要如此理解。因为如果圣言是通过变形和实体的变化成为肉身，那么立刻就会得出：耶稣必定是一个由两个实体——肉身和圣神——复合而成的实体，一种混合物，如同\Emph{\OriginalTerm{金银合金}{electrum}}，由金和银组成；于是它既不是金（即神）也不是银（即肉身）——一个被另一个改变，产生了第三个实体。因此，按此方式，耶稣不能是天主，因为祂不再是成了肉身的圣言；也不能是降生成人的真人，因为祂并非真正的肉身，而圣言所成为的正是肉身。因此，被两者复合而成，祂实际上什么都不是；毋宁说祂是某个第三者实体，与两者截然不同。但事实是，我们发现祂被明确展示为既是天主又是人；我们所引用的圣咏暗示（关于肉身）：\Quote{天主在其中成为人，祂因此因圣父的旨意建立了它，}——无论从哪方面看，作为天主子和人子，既是天主又是人，无疑按照每个实体在其自身特殊属性中有所区别，因为圣言无非是天主，肉身无非是人。宗徒也同样教导关于祂的两个实体，说：\Quote{祂按肉身是达味的后裔；}\BibleRef{罗 1:3} 在这话中，祂将是人子和人。\Quote{祂按圣神被立为天主子；}在这话中，祂将是天主，并且是圣言——天主子。我们清楚地看到这双重状态，并不混乱，而是结合在同一个位格内——耶稣，天主和人。关于基督，我暂且推迟我要说的话。（我在此指出），每一性体的特性被如此完整地保存，以至于一方面圣神在耶稣内做了适合祂的一切事，如奇迹、大能和异事；另一方面，肉身展示了属于它的情感。它在魔鬼的诱惑下饿了，与撒玛黎雅妇人一起渴了，为\Person{拉匝禄}哭了，忧愁至死，最后确实死了。然而，如果它只是一个\Emph{\OriginalTerm{第三者}{tertium quid}}，一些由两个实体形成的复合本质，就像我们提到的\Emph{\OriginalTerm{金银合金}{electrum}}，那么就不会有任何明显证据表明任一性体。但通过功能的转移，圣神会做该由肉身做的事，肉身也会做该由圣神做的事；或者做既不适合肉身也不适合圣神的事，而是混乱地具有某种第三者特性。不仅如此，在这种假设下，要么圣言经历了死亡，要么肉身没有死——如果圣言被变成了肉身；因为要么肉身是不死的，要么圣言是会死的。然而，既然两个实体各自按其特性分别行动，它们就必然各自有其运作和结果。那么，与\Person{尼苛德摩}一起学习：\Quote{那由肉身生的就是肉身，那由圣神生的就是圣神。}\BibleRef{若 3:6} 肉身不会成为圣神，圣神也不会成为肉身。在一个\Emph{位格}内，它们无疑能够共存。耶稣由它们构成——人，出于肉身；出于圣神，天主——天使称祂为\Quote{天主子，}\BibleRef{路 1:35} 这是按祂是圣神的那一性体，而给肉身保留了\Quote{人子}的称号。同样，宗徒又称祂为\Quote{天主与人之间的中保，}\BibleRef{弟前 2:5} 从而确认了祂对两个实体的参与。现在，为了结束此事，你们这些把天主子解释为肉身的人，请告诉我们人子是什么？我想知道，祂会是神吗？但你们坚持认为圣父自己就是神，理由是\Quote{天主是神，}好像我们没读到还有\Quote{天主之神}似的；同样，我们发现\Quote{圣言就是天主，}但也有\Quote{天主的圣言}。\par

% structure: p0057 source=h2 target=section reason=relative-heading-rank; base=h2

\section{第二十八章. 基督不是圣父，如\Person{普拉克色阿斯}所说。此意见的不一致性，不亚于其荒谬性，被揭露。耶稣基督的真道，根据\Person{圣保禄}，他与其他神圣作者一致。}

所以，最愚蠢的异端者啊，你使\Emph{基督}成为圣父，却从未考虑这个名称的实际力量，如果基督是一个名字，而不是一个姓氏或称号；因为它意味着\OriginalTerm{傅油者}{Anointed}。但\Quote{傅油者}并不比\Quote{着衣者}或\Quote{穿鞋者}更是一个专名；它只是名字的一个附属品。现在假设，某种方式下，耶稣也被称为\Emph{着衣者}（Vestitus，\Emph{Clothed}），正如祂因祂傅油的奥秘而实际被称为基督，你是否会同样说耶稣是天主子，同时又假设\Emph{着衣者}是圣父？那么现在，关于基督，如果基督是圣父，那么圣父就是一位傅油者，并且当然从另一位领受傅油。否则，如果祂是从自己领受的，那么你必须向我们证明。但我们从\WorkTitle{宗徒大事录}中教会的呼求里学不到这样的事实：\BibleQuote{确实，主啊，反对祢的圣童耶稣，祢所傅油的，黑落德和般雀比拉多与外邦人及以色列人民聚集在一起。}\BibleRef{宗 4:27}这些事便同时证明耶稣是天主子，并且作为子，祂由圣父傅油。因此，基督必然与那位由圣父傅油的耶稣是同一的，而不是圣父，是圣父傅了子。\Person{伯多禄}的话也有同样的效力：\BibleQuote{以色列全家应确切知道，天主已使你们所钉死的这位耶稣成为主与基督，也就是傅油者。}\BibleRef{宗 2:36}此外，\Person{若望}谴责那\Quote{否认耶稣是基督}的人为\Quote{说谎者}；另一方面他宣告\Quote{凡相信耶稣是基督的，都是由天主生的。}因此他也劝我们信祂（\Emph{圣父的}）子耶稣基督的名，好使\BibleQuote{我们的共融与圣父并与祂的子耶稣基督相通。}\BibleRef{若一 1:3}\Person{保禄}也同样处处提到\Quote{天主圣父和我们的主耶稣基督。}当他写信给罗马人时，他藉我们的主耶稣基督感谢天主。\BibleRef{罗 1:8}给迦拉达人时，他宣告自己是\BibleQuote{非由人，亦非藉着人，而是藉着耶稣基督和天主圣父成为宗徒。}\BibleRef{迦 1:1}你确实拥有他所有的著作，都清楚地见证同一事实，并表明两位——天主圣父和我们的主耶稣基督，圣父之子。（它们也见证）耶稣本身就是基督，并且无论用哪一个称号，都是天主子。因为正是基于同样的权利，两个名字都属于同一位格，即天主子，任何一个名字单独出现也属于同一位格。因此，无论是单独出现耶稣这个名字，基督也被理解为同一位，因为耶稣是傅油者；或者如果只给出基督这个名字，那么耶稣就等同于祂，因为傅油者就是耶稣。现在，在这两个名字\Quote{耶稣基督}中，前者是专名，由天使赋予祂；后者只是一个附属品，从祂的傅油可以称呼祂——从而暗示一个条件：基督必是子，而非圣父。何其盲目，的确，那人不察觉，若他将基督这个名字归于圣父，则藉基督之名暗示了另一位天主！因为如果基督是天主圣父，当祂说：\BibleQuote{我升到我的父和你们的父那里，到我的天主和你们的天主那里}\BibleRef{若 20:17}时，祂当然足以清楚地表明，在祂之上有另一位父和另一位天主。如果，再者，圣父是基督，那么祂必定是另一位存在者，祂\Quote{加强雷霆，创造风，并向人宣告祂的基督。}并且如果\Quote{地上的君王起来，首领聚集在一起反对上主和反对祂的基督，}那么上主必定是另一位存在者，反对祂的基督，君王和首领聚集在一起。并且如果，引用另一处经文：\Quote{上主对我主基督说，}那位对基督的父说话的上主必定是另一位存在者。再者，当宗徒在他的书信中祈祷：\BibleQuote{愿我们主耶稣基督的天主赐给你们智慧和知识的神恩}\BibleRef{弗 1:17}时，祂——耶稣基督的天主、属灵恩赐的赐予者——必定是另一位（不同于基督的）。并且一言以蔽之，免得我们遍寻各处经文，那位\BibleQuote{从死人中复活了基督，也将复活我们必死的身体}\BibleRef{罗 8:11}，作为使活者，必定不同于死了的圣父，甚至不同于被复活的圣父——如果死去的基督是圣父的话。\par

% structure: p0059 source=h2 target=section reason=relative-heading-rank; base=h2

\section{第二十九章. 是基督死了。圣父不能受苦，无论是独自还是与他人一起。从\Person{普拉克色阿斯}的前提产生的亵渎结论。}

住口！对这亵渎住口。让我们满足于这样说：基督死了，圣父之子；并且\Emph{让这足够了}，因为圣经已经如此告诉我们。因为就连宗徒，在他那并非不感到其沉重的宣告中——即\Quote{基督死了}——他立刻加上\Quote{照圣经所记载的}，\BibleRef{格前 15:3}以便以圣经的权威缓和这话的严厉，从而除去读者的反感。现在，既然在基督内被宣称有两种实体——即天主性和人性——显然由此得出结论：天主性是不死不灭的，而人性是有死的；那么，他在何种意义上宣称\Quote{基督死了}就是明显的——即在他曾是肉性和人、及人子，而非圣神和圣言及天主之子的意义上。简言之，既然他说是\Emph{基督}（即受傅者）死了，他就向我们指明，死去的是那受傅的本性，即肉性。很好，你说；既然我们这一方面用你们关于圣子所用的相同措辞来肯定我们的道理，我们就不是对主天主犯下亵渎之罪，因为我们不主张祂是按天主性死了，而只是按人性。不，你是在亵渎；因为你不仅说圣父死了，而且说祂是死在十字架上。因为\Quote{凡挂在树上的，是可咒骂的}，\BibleRef{迦 3:13}——这咒骂，按照法律，是适用于圣子的（因为\Quote{基督为我们成了咒骂}，但绝不是圣父）；然而，既然你把基督变成了圣父，你就可被指为对圣父的亵渎。但当我们断言基督被钉在十字架上时，我们并非以咒骂来毁谤祂；我们只是重申法律所宣布的咒骂：\BibleRef{申 21:23}而且宗徒在说同样的话时也确实没有亵渎。\BibleRef{迦 3:13}此外，既然将那恰如其分地属于某主体的东西归于该主体并非亵渎；那么，相反，将不适于该主体的东西归于它则是亵渎。按照这一原则，圣父也没有与圣子同受苦难。\Emph{异端者}害怕直接冒犯对圣父的亵渎，期望通过这一权宜之计来减轻它：他们承认我们至此，即圣父和圣子是两位；又加上说，既然确实是圣子受苦，圣父只是祂的同受苦难者。但他们在这一自负中多么荒谬！因为\Quote{同受苦难}是什么意思，不就是与另一个人一起承受苦难？现在如果圣父是不能受苦的，那么祂就不能与另一个人一同受苦；否则，如果祂能与人同受苦，那么祂当然也能受苦。事实上，你通过这恐惧的遁词一点也没有让步。你不敢说祂能受苦，却使祂能够同受苦。再者，圣父不能同受苦，正如圣子按祂作为天主的存在条件也不能受苦一样。好了，但圣子怎能受苦，如果圣父不与祂同受苦？\Emph{我的回答是}：圣父是与圣子分开的，虽然并非与作为天主的祂分开。因为即使一条河被淤泥和泥土玷污，虽然它从同一本性的源泉流出，且不与源泉分离，但影响河流的损害却不会到达源泉；并且虽然河流中的水是源泉的水在受苦，但由于它不是在源泉处受损害，而只是在河中，所以源泉什么也没有受，只有从源泉流出的河流受苦。同样，天主圣神，无论祂可能在圣子内承受何种苦难，然而，既然祂不能在圣父（\Emph{天主性的泉源}）内受苦，而只能在圣子内受苦，祂显然不能作为圣父受苦。但对我而言，天主圣神作为天主圣神什么也没有受就足够了，因为祂所受的一切都是在圣子内受的。圣父与\Emph{在肉身中的圣子}同受苦完全是另一回事。这也已经被我们处理过。没有人会否认这一点，因为甚至我们自己也不能为天主受苦，除非天主圣神在我们内，祂也藉我们的器皿说出关于我们行为和受苦的一切；然而，祂并非在祂自己内受苦，只是赐给我们受苦的能力和容量。\par

% structure: p0061 source=h2 target=section reason=relative-heading-rank; base=h2

\section{第三十章 圣子如何在十字架上被圣父舍弃；其真正意义对\Person{普拉克塞阿斯}是致命的；同样，基督的复活、升天、坐在圣父的右边，以及圣神的派遣}

然而，如果你坚持进一步推行你的观点，我将找到更严厉的方式来回应你，并以主亲自的呼喊来迎接你，用这个问题来质问你：你究竟在探究和推理\Emph{那个}？你在祂受难时听见祂呼喊：\BibleQuote{我的天主，我的天主，祢为什么舍弃了我？}\BibleRef{玛 27:46} 那么，要么是圣子受苦，被圣父\Quote{舍弃}，而圣父因此没有受苦，因为祂舍弃了圣子；要么，如果是圣父受了苦，那么祂是向哪位天主呼喊呢？但这血肉和灵魂的声音，即人的声音——不是圣言和圣神的声音，即不是天主的声音——发出，是为了证明天主的不可受难性，祂\Quote{舍弃}了祂的圣子，将祂的人性交出经历死亡。这位宗徒也领悟了这真理，他如此写道：\BibleQuote{如果圣父连自己的圣子都没有怜惜}\BibleRef{罗 8:32}。依撒意亚在此之前也同样领悟了，他宣告：\Quote{主为我们的过犯交付了祂。} 就这样，祂\Quote{舍弃}了祂，即\Emph{没有怜惜}祂；\Quote{舍弃}了祂，即\Emph{交出}祂。在其他一切方面，圣父并没有舍弃圣子，因为圣子将祂的灵魂交在了圣父手中。\BibleRef{路 23:46} 确实，在这样交托之后，祂立刻死了；既然圣神仍与肉身同在，肉身就不能经历完全的死亡，即\Emph{腐朽与败坏}。因此，对圣子来说，死就等于被圣父舍弃。然后，圣子死而复活，依照圣经。\BibleRef{格前 15:3} 也是圣子升到高天，\BibleRef{若 3:13} 又下到地下深处。\BibleRef{弗 4:9} \BibleQuote{祂坐在圣父的右边}——而非圣父坐在自己的右边。\Person{斯德望}在被人用石头砸死殉道时，看见祂仍然坐在天主的右边，\BibleRef{宗 7:55} 祂将继续坐在那里，直到圣父使祂的仇敌成为祂的脚凳。祂将乘着天上的云彩再来，正如祂升天时所显现的那样。同时，祂从圣父那里领受了所应许的恩赐，并将它倾注下来，即圣神——天主性中的第三位，神圣威严的第三级；\Emph{独一天主神权}的宣告者，同时也是\Emph{经世}的解释者，向每一位听见并接受新预言话语的人；又是\BibleQuote{引入一切真理的那位}，\BibleRef{若 16:13} 即存在于圣父、圣子和圣神中的那位，按照基督道理的奥秘。\par

% structure: p0063 source=h2 target=section reason=relative-heading-rank; base=h2

\section{第三十一章 \Person{普拉克塞亚}异端的倒退性质。\OriginalTerm{天主圣三}{Blessed Trinity}的教义构成犹太教与基督教之间的重大区别。}

但是，（你们的这个教义）与犹太信仰相似，其本质就是：相信一个天主，以至于拒绝将圣子与祂并列，以及在圣子之后的圣神。现在，如果没有这个\Emph{你们正试图打破的}区别，我们和他们之间还有什么不同呢？如果从此以后，圣父、圣子和圣神不被相信是三位，并且是唯一的天主，那么福音（它是新盟约的本质）还有什么必要规定（正如它所规定的）法律和先知直到若翰\Emph{洗者}为止呢？天主乐意以这样的方式与人更新祂的盟约，即祂的合一可以以一种新的方式，通过圣子和圣神被相信，好使天主现在能被公开地认识，在祂固有的名号和位格中，而这位天主在古代虽然通过圣子和圣神被宣告，却不曾被清楚地理解。那么，让那些\Quote{否认圣父和圣子的假基督}离开吧。因为他们说圣父与圣子相同，就否认了圣父；他们假设圣子与圣父相同，就否认了圣子，他们将不属于\Emph{他们的}事物归于\Emph{他们}，并夺走属于\Emph{他们的}事物。但\BibleQuote{凡承认（耶稣）基督是天主子的}（而非圣父），\BibleQuote{天主就住在他内，他也住在天主内。}\BibleRef{若一 4:15} 我们不相信天主为祂的圣子所作的见证。\BibleQuote{没有圣子的，就没有生命。}\BibleRef{若一 5:12} 而凡相信祂不是圣子以外的任何其他存在的，就没有圣子。\par

\SourceRecord{Translated by Peter Holmes.；Ante-Nicene Fathers；Vol. 3.；Edited by Alexander Roberts, James Donaldson, and A. Cleveland Coxe.；Buffalo, NY: Christian Literature Publishing Co.,；1885.；<http://www.newadvent.org/fathers/0317.htm>.}
